% \iffalse meta-comment
%
% Copyright (C) 2017- by Yanshuo Chu <yanshuoc@gmail.com>
%
% This file may be distributed and/or modified under the
% conditions of the LaTeX Project Public License, either version 1.3a
% of this license or (at your option) any later version.
% The latest version of this license is in:
%
% http://www.latex-project.org/lppl.txt
%
% and version 1.3a or later is part of all distributions of LaTeX
% version 2004/10/01 or later.
%
% \fi
%
% ^^A ======================================================================
% ^^A hithesis 用户手册
% ^^A
% ^^A 只放面向用户的说明（模板介绍 / 安装 / 使用说明），不含任何实现代码，
% ^^A 也不参与 docstrip：hithesis.ins 从不读这个文件。
% ^^A 由 hithesis.dtx 的 driver 通过 \DocInput 引入，渲染成手册的前半部分。
% ^^A
% ^^A 手册由 hithesis 自己排成一本博士学位论文（driver 里的 \documentclass）：
% ^^A 这里先给封面的元信息与中英文摘要，再起前置与正文，三章正文接在后面。
% ^^A 附录（实现）、修改记录、索引与致谢由 driver 接在这个文件之后。
% ^^A ======================================================================
%
% \GetFileInfo{\jobname.dtx}
%
% \changes{v3.2a}{2026/09/20}{手册排成一本博士学位论文：封面、摘要、目录、章节、附录、
%   索引与致谢走类自己的版面，正文三章介绍、安装与使用，实现进附录，修改记录与索引
%   排在后置；免责声明更新，维护者名单挪到致谢}
% \changes{v3.2a}{2026/09/20}{安装与生成一章按 v3.2a 重写：\hologo{TeX} Live 版本口径、
%   推荐引擎改成 \hologo{LuaLaTeX}（断行与行高照 Word 的那一套只有它有）、
%   \file{hithesis.dtx} 自解压、\texttt{make~manual} 改名；加一节列出相对 v3.1e
%   的破坏性变更（类名、\option{degree-level}、\option{stage}、\cs{pozhehao}、
%   示例宏包 \file{hithesis.sty}、字体选项、脚注、中文斜体、双语题注）}
% \changes{v3.2a}{2026/09/13}{使用说明一章按 v3.2a 的接口重写：\cs{hitsetup} 的族表、
%   条件参数与先后规矩、\option{struct} 族与骨架宏、\option{defense} 族与答辩决议页、
%   废弃的写法、\option{zh}／\option{en} 分组与斜杠路径、日期字段、符号表与缩略语表、
%   \pkg{tabularray} 表格与长表格、\pkg{subcaption} 子图、\cs{figurenote}、
%   \env{eqdenote}、\pkg{thmtools} 定理、\cs{hitenter}、成果分组、\option{term} 族的
%   键名构成与 \file{hithesis.cfg}}
% \changes{v3.2a}{2026/09/21}{手册与实现说明里的表一律改成类的 \env{longtblr}：三线表、
%   带表题与表号、长的列用 \texttt{X} 列折行、跨页自动续表头；原先是 \env{longtable}
%   与 \env{tabular} 的 \texttt{l} 列，长内容压出版心}
%
% \hitsetup{ info = {
%   secret-level = {公开},
%   zh = {
%     title = {\hithesis：哈尔滨工业大学}{学位论文模板},
%     author = {孙忠豪\quad 石佳乐\quad 李云鹏\quad 支卓\quad 吴俊豪},
%     supervisor = {初砚硕},
%     affiliation = {GitHub hithesis/hithesis},
%     major = {文档排版},
%     discipline = {工学},
%     keywords = {{学位论文模板}, {哈尔滨工业大学}, {\hologo{LaTeX}}, {开题报告}, {中期报告}},
%   },
%   en = {
%     title = {\hithesis: The Thesis Template of Harbin Institute of Technology},
%     author = {Sun Zhonghao, Shi Jiale, Li Yunpeng, Zhi Zhuo, Wu Junhao},
%     supervisor = {Chu Yanshuo},
%     affiliation = {GitHub hithesis/hithesis},
%     major = {Document Typesetting},
%     discipline = {Engineering},
%     keywords = {thesis template, Harbin Institute of Technology, LaTeX, proposal, interim report},
%   },
% } }
% ^^A 封面与内封上的字样照门派的叫法改：历任维护者是掌门，初砚硕是太上长老，
% ^^A 手册是掌门秘籍；页眉不跟封面，印手册自己的名字。手册单面看，不跳右手页。
% \hitsetup{
%   layout = { openright = false },
%   term = {
%     base-degree-level-zh = {掌门},
%     base-document-type   = {秘籍},
%     header = { document-type = {} },
%     zh = {
%       header = { university = {\hithesis\ 用户手册}, degree-level = {} },
%       titlepage = {
%         author = {掌门}, supervisor = {太上长老},
%         degree-applied = {修为}, degree-of-discipline = {工学博士},
%         major = {所修功法}, affiliation = {山门},
%         date = {出关日期}, institution = {师门},
%       },
%     },
%     en = { titlepage = {
%       author = {Sect Masters}, supervisor = {Grand Elder},
%       degree-applied = {Cultivation}, major = {Martial Art},
%       affiliation = {Headquarters}, date = {Date of Emergence},
%       institution = {Alma Mater},
%     } },
%     titlepage-degree-sentence = {Secret Manual of the Sect Masters},
%   },
% }
%
% \begin{abstract-zh}
% \hithesis\ 是哈尔滨工业大学一校三区的 \hologo{LaTeX} 论文模板，支持本科、
% 硕士、博士终稿，开题与中期报告，英文论文，以及博士后出站报告。文档类负责
% 封面、摘要、目录、正文和后置材料的版式与顺序，作者填写元信息并组织正文即可。
%
% 本手册也由 \hithesis\ 排版。正文介绍安装、配置和常用命令；实现代码放在附录，
% 修改记录与命令索引接在其后。
% \end{abstract-zh}
%
% \begin{abstract-en}
% \hithesis\ is a \hologo{LaTeX} class for theses, dissertations, proposals,
% interim reports and postdoctoral reports of Harbin Institute of Technology
% on all three campuses. The class selects the required layout from the campus,
% degree and document stage, and arranges the front matter, body and back matter.
%
% This manual is typeset with \hithesis{} itself. Its main text covers installation,
% configuration and user commands; the implementation, change history and command
% index follow in the appendix and back matter.
% \end{abstract-en}
%
% \hitfrontmatter
% \hitmainmatter
%
% \chapter{模板介绍}
% \hithesis\ 是哈尔滨工业大学学位论文与报告的文档类。校区、学位、正文语言和
% 交付物在类选项里声明；标题、作者等元信息通过 \cs{hitsetup} 填写。先从
% \file{examples/demo/} 复制一份示例，通常比从空文件开始省事。
%
% 有疑问时先查对应章节和示例；仍不能定位的问题，按第~\ref{sec:howtoask}~节
% 准备最小示例与日志。源码维护约定见仓库根目录的 \file{CONTRIBUTING.md}。
%
% \section{本文使用的记号}
% 宏包名和文档类用无衬线体，如 \pkg{xeCJK} 与 \cls{hithesis}；命令、环境和
% 文件名用等宽体，如 \cs{hitsetup}、\env{abstract} 与 \file{final.tex}。
% 尖括号里的文字表示需要替换的参数，输入时不写尖括号。蓝色侧边线标出
% \hologo{LaTeX} 代码，橙色侧边线标出命令行。
%
% \note[注意：]{模板只替你处理已经实现的版式。擅自改动结构、跳过必填字段或混用
% 旧接口，输出自然可能偏离学校规范。}
%
% \section{免责声明}
% 本模板依据哈尔滨工业大学教务处和研究生院公开的写作指南与书写范例实现，不是
% 学校官方模板。规范与范例相互矛盾时，项目只能选择一套可复核的排法；评审老师
% 仍可能提出不同要求。使用前请确认所在院系当年的通知。
%
% 源码按 \hologo{LaTeX} Project Public License 发布。修改和再发布时应遵守许可证。
%
% \chapter{安装}
% \label{sec:installation}
% 模板支持 \hologo{TeX} Live 2022 及以上版本，排版基准是 2026，推荐使用基准版。
% 2024、2025 与基准版的主要页面断点已经对齐；2022、2023 受旧版
% \pkg{tabularray} 等上游宏包影响，个别长表间距会有差别。各版本都应能正常编译。
%
% 稳定版可从 CTAN 或 GitHub Releases 获取；\texttt{dev} 分支是开发版。论文已经
% 进入定稿阶段时，升级前先保留旧环境和 PDF，再完整比较一遍输出。
%
% 项目地址：
% \begin{itemize}
% \item \href{https://github.com/hithesis/hithesis}{GitHub 仓库}
% \item \href{https://github.com/hithesis/hithesis/releases}{发行版}
% \item \href{https://ctan.org/pkg/hithesis}{CTAN}
% \end{itemize}
%
% \section{模板的组成}
% 下表列出了 \hithesis\ 的主要文件及其功能介绍：
%
% \begin{longtblr}[caption={模板的主要文件}, label={tab:files}]
%   {colspec={l X[2,l]}, width=\linewidth}
% \toprule
% 文件（夹） & 功能描述 \\ \midrule
% \file{hithesis.ins} & \textsc{DocStrip} 驱动文件，由 \file{hithesis.dtx} 生成 \\
% \file{hithesis.dtx} & \textsc{DocStrip} 主源文件（开发用）\\
% \file{hit-*.dtx} & 类的各个模块源文件（开发用）\\
% \file{hit-eps.dtx} & EPS 图像数据（校徽、封面、示例插图等，从主源剥离）\\
% \file{dtx-style.sty} & \textsc{DocStrip} 样式文件（开发用）\\\midrule
% \file{hithesis.cls} & 类文件，终稿与开题中期都是它，由 \option{stage} 分 \\
% \file{hithesisbook.cls} & 旧类名，已移除，写了会报错说明改法 \\
% \file{hithesisart.cls} & 旧类名，已移除，写了会报错说明改法 \\
% \file{hithesisartplus.cls} & 旧类名，已移除，写了会报错说明改法 \\
% \file{hithesis.bst} & 参考文献样式文件 \\\midrule
% \file{hithesis.ist} & 索引样式文件\\
% \file{hithesis.cfg} & 固定词汇与各项默认值，两档共用一份，见第 \ref{sec:cfg} 节 \\
% \file{hithesis.sty} & 示例宏包，示例文档自己 \cs{usepackage}，见第 \ref{sec:example-package} 节 \\\midrule
% \file{examples/demo/} & 模板示例文件夹，学位论文与开题中期都在这里 \\
% \texttt{├─}~\file{final.tex} & 学位论文主文件，中英文两版共用 \\
% \texttt{├─}~\file{report.tex} & 开题报告与中期报告主文件 \\
% \texttt{├─}~\file{info.tex} & 两份主文件共用的元信息 \\
% \texttt{├─}~\file{thesis/} & 学位论文的各部件 \\
% \texttt{└─}~\file{report/} & 两种报告的各部件 \\\midrule
% \file{Makefile} & Makefile\\
% \file{latexmkrc} & latexmk 配置文件 \\
% \file{README.md} & Readme\\
% \file{hithesis.pdf} & 用户手册（本文档）\\
% \bottomrule
% \end{longtblr}
%
% 生成后的 \file{examples/demo/} 可以整体复制到论文目录。平时写论文只需要
% \file{hithesis.cls}、\file{hithesis.cfg}、参考文献与索引样式、示例宏包和自己的
% 正文；\file{src/} 下的 \file{.dtx} 是开发源码。
%
% \section{代码组织}
% \label{sec:modules}
% 类的实现按职责拆在 \file{src/engine/}、\file{src/config/}、\file{src/flow/}
% 和 \file{src/pages/} 下，生成时拼成一个 \file{hithesis.cls}。这部分只与维护者
% 有关；模块顺序、公开接口和检查命令见 \file{CONTRIBUTING.md}。
%
% \section{expl3 与空格}
% \label{sec:expl3-space}
% v3.2a 起类文件从头到尾都在 \pkg{expl3} 的 catcode régime 里，由开头的
% \cs{ProvidesExplClass} 打开，中间不再有成对的 \cs{ExplSyntaxOn} 与
% \cs{ExplSyntaxOff}。这套 catcode 下空格的含义和平时不一样，改代码时有三件事要留神。
%
% \textbf{字面空格会被吞。} \texttt{text=\{150true mm\}} 照抄进去会变成
% \texttt{150truemm}，要空格得写 \texttt{\~{}}。控制词后面的空格是名字终结符，
% 本来就会被吸收，不用管。行末的换行同样被忽略，多行的文本值里原本靠换行产生的
% 那个空格要补 \texttt{\~{}}，原本靠空行产生的 \cs{par} 要显式写出来。
%
% \textbf{\texttt{\~{}} 的含义反过来了。} 平时它是活动字符，表示不断行空格；
% \pkg{expl3} 下它是普通空格。原有代码里的 \texttt{\~{}} 照抄进来，不断行空格
% 就退化成了普通空格，要保持原意得写 \cs{nobreakspace}\{\}。
%
% \textbf{控制空格不在此列}，它在两种语法下是同一个控制
% 序列，照写即可。
%
% 另有两处例外，必须退回 2e 的 catcode：走裸 \cs{input} 的加载器（如
% \cs{usetikzlibrary}），以及 \pkg{etoolbox} 的 \cs{patchcmd} 一族——它们会把
% 文件或宏体按当前 catcode 重扫，字面空格照样被吞。v3.2a 已经把这两类都消掉了，
% 类文件里一处 \cs{ExplSyntaxOn} 都没有。
%
% \section{从源码生成模板}
% \label{sec:generate-cls}
% 从 CTAN 安装的用户不需要执行本节。使用 git 仓库时，在仓库根目录运行：
% \begin{shell}
% $ make cls
% \end{shell}
% 该目标生成类文件、配置、参考文献与索引样式，并复制到示例目录。
%
% 没有 \texttt{make} 时可直接解包：
% \begin{shell}
% $ xetex src/hithesis.dtx
% \end{shell}
% CTAN 源码包也可运行 \texttt{latex hithesis.ins}。仓库源码含中文，直接跑
% \file{src/hithesis.dtx} 时要用 \texttt{xetex}，plain \texttt{tex} 会把中文写成
% 转义序列。
%
% 用户手册不是使用模板的前置条件。需要本地生成时运行：
% \begin{shell}
% $ make manual
% \end{shell}
%
% \section{生成论文}
% \label{sec:generate-thesis}
% 本节介绍几种常见的生成论文的方法。用户可根据自己的情况选择。
%
% \subsection{用哪个引擎}
% \label{sec:engine}
% 模板认 \texttt{xelatex} 与 \texttt{lualatex} 两个引擎，\texttt{pdflatex} 与
% \texttt{latex} 不行——那两个不认 OpenType 字体，排不出本模板要的中文。写错引擎
% 会在第一时间报 \texttt{Please use: lualatex (or xelatex)}，不必等一屏字体错误。
%
% \textbf{用 \texttt{lualatex}。} 2026 年 9 月起模板的断行与行高照 Word 排：每一行
% 按 Word 的贪心算法首次适配、装不下就压标点，行盒的高深按行里实际的字体逐行算。
% 这两件事都在 Lua 里做（\file{src/docxlike/docxlike-linebreak.dtx}），只有 \hologo{LuaLaTeX}
% 有这条路，与 Word 逐行逐页的对拍也都是拿它跑的。示例的魔法注释、\file{Makefile}
% 与 \file{latexmkrc} 默认都是它。
%
% \hologo{XeLaTeX} 仍能编，但退化成 \hologo{TeX} 自己的断行加 \pkg{xeCJK} 的标点
% 处理，行高按段落样式那副字撑杆、不逐行看内容，与 Word 逐行相同就谈不上了；
% 另有两处功能上的出入，在 \file{src/engine/hit-engine.dtx} 里写了原委：
% \begin{itemize}
%   \item \hologo{LuaLaTeX} 下 \cs{CJKunderanyline} 与 \cs{CJKunderanysymbol} 没有，
%     那边顶替 \pkg{xeCJKfntef} 的是 \pkg{ulem}，它没有对得上的东西。其余
%     \cs{CJKunderline}、\cs{CJKunderdblline}、\cs{CJKunderwave}、\cs{CJKsout}、
%     \cs{CJKxout}、\cs{CJKunderdot} 都在。
%   \item \cs{quan} 画圈时（字体没有现成的带圈数字才会画）圈里的数字取自西文字体，
%     两个引擎粗细不同。常见的几档中文字体都自带
%     \texttt{①}--\texttt{⑳}，多数文档走不到这一支。
% \end{itemize}
% 要交的稿子中途别换引擎。下面几小节的命令要用 \hologo{XeLaTeX} 时把
% \texttt{lualatex} 换成 \texttt{xelatex}，\texttt{latexmk} 那条换成 \texttt{-xelatex}。
%
% \subsection{手工编译}
% 下面以示例的 \file{final.tex} 和传统 BibTeX 后端为例：
% \begin{shell}
% $ lualatex final
% $ bibtex final
% $ lualatex final
% $ lualatex final
% $ splitindex final -- -s hithesis.ist
% $ lualatex final
% \end{shell}
% 不生成主题索引时可省去最后两步。选择 biber 后端时把 \texttt{bibtex final}
% 换成 \texttt{biber final}。日志仍提示交叉引用变化时继续运行 LuaLaTeX。
%
% \subsection{latexmk}
% 示例目录带有 \file{latexmkrc}，会按文档实际用到的参考文献和索引工具反复编译，
% 直到交叉引用稳定：
% \begin{shell}
% $ cd examples/demo
% $ latexmk final.tex
% \end{shell}
% 开题或中期报告把文件名换成 \file{report.tex}。
%
% \subsection{make}
% 仓库根目录的 \file{Makefile} 负责生成类文件和手册：
% \begin{shell}
% $ make cls
% $ make manual
% \end{shell}
% \file{examples/demo/Makefile} 负责编译示例：
% \begin{shell}
% $ make final
% $ make report
% \end{shell}
% 两份 Makefile 默认都调用 LuaLaTeX。需要改引擎时，应同时修改 \texttt{METHOD}
% 与 \texttt{LATEXMKOPTS}，不要让直接编译和 \texttt{latexmk} 走不同引擎。
%
% \section{升级}
% \label{sec:upgrade}
% 通过 TeX 发行版安装的用户，用发行版的包管理器更新。使用仓库副本的用户，拉取
% 新代码后重新运行 \texttt{make cls}，再完整编译自己的论文。不要只复制一个
% \file{hithesis.cls}：配置、参考文献样式、索引样式和 Lua 断行文件也可能同时变化。
%
% 跨版本升级先读下一节的接口变更。定稿阶段还应保留升级前的 PDF，逐页核对版式。
%
% \section{从 v3.1e 升上来要注意的}
% \label{sec:breaking}
% v3.2a 把底层全部迁到 \pkg{expl3}，对外接口大体不变，但有几处确实变了。
% 下面按“会不会编不过”排序。
%
% \subsection{编不过的}
%
% \begin{itemize}
% \item \textbf{删掉了 \cs{dachu}、\cs{banxiaosi}、\cs{dawu} 三个字号。}
%   这三个不在通行的中文字号里，手册的字号对照表从来没列过。用到的改成相邻的
%   标准字号：\cs{dachu}（58bp）改 \cs{chuhao}，\cs{banxiaosi}（13bp）改
%   \cs{sihao} 或 \cs{xiaosi}，\cs{dawu}（11bp）改 \cs{xiaosi} 或 \cs{wuhao}。
% \item \textbf{删掉了 \option{pifootnote} 选项与 \pkg{pifont}。}
%   带圈脚注号现在统一由 \cs{quan} 排，先取当前宋体自带的带圈字符，超出字体可用
%   范围才画圈。要 Dingbats 那套 ➀➁➂ 的自己 \cs{usepackage}\marg{pifont}。
% \item \textbf{不再预载 \pkg{tikz} 的 \texttt{calc} 库。}
%   自己的插图里写了 \texttt{(\$…\$)} 坐标运算的，在导言区补一句
%   \cs{usetikzlibrary}\marg{calc}。
% \item \textbf{删掉了 \cs{pozhehao}。}
%   它只是个排“——”的宏，占了一个公开名字却什么也没多做。正文里用到的直接打
%   “——”两个字符。

% \item \textbf{最低 \hologo{TeX}~Live 版本提到 2022。}
% \end{itemize}
%
% \subsection{能编过但要留意的}
%
% \begin{itemize}
% \item \textbf{目录里自己改过的 \file{hithesis.sty} 照旧优先。}
%   \cs{usepackage}\marg{hithesis} 这一行不用动，加载到的还是自己那份，改过的内容
%   都还在。
%
%   v3.1e 那份里 \cs{RequirePackage}\marg{bm} 没有条件，而 \pkg{bm} 跟 v3.2a 默认的
%   \pkg{unicode-math} 不兼容，公式里一出现希腊字母就报
%   \texttt{Extended mathchar used as mathchar}，完整示例实测 70 个。类里挂了钩子：
%   \pkg{bm} 真被装上时，\cs{bm} 与 \cs{boldsymbol} 会接回 \pkg{unicode-math} 自己的
%   \cs{symbf}，日志里说明一句，文档不用改。确实要用 \pkg{bm} 本身的加
%   \option{math-font=newtx} 走 8 位数学字体。
% \item \textbf{\option{type} 改名 \option{degree-level}。}
%   v3.2a 加了 \option{stage}（选交付物是终稿还是开题中期），于是 \option{type}
%   与 \option{stage} 两个名字都读作“种类”，一个管学位一个管交付物，谁也说不清
%   哪个是哪个。学位那个改叫 \option{degree-level}：
%
%   \begin{latex}
%   \documentclass[degree-level=doctor,campus=harbin]{hithesis}
%   \end{latex}
%
%   旧名照旧生效，编译时提示一次。它是加载期选项，改 \cs{documentclass} 那一行就行，
%   不要挪进 \cs{hitsetup}。旧名在 v4（预计 2027 年发布）移除。
% \item \textbf{开题与中期用 \option{stage} 选。}
%   v3.1e 是另一个类名（\file{hithesisart}），v3.2a 起是同一个类加一个选项：
%
%   \begin{latex}
%   \documentclass[stage=proposal,degree-level=doctor,campus=harbin]{hithesis}
%   \end{latex}
%
%   两个旧取值 \option{stage}\texttt{=opening} 与 \option{stage}\texttt{=midterm}
%   保留，按 \texttt{proposal} 与 \texttt{interim} 处理，编译时提示一次，
%   在 v4（预计 2027 年发布）移除。
% \item \textbf{文档类改名了，三个旧名已移除。}
%   现在统一写 \file{hithesis}，交付物由 \option{stage} 选。改法只是把
%   \cs{documentclass} 那一行的类名换掉，文档其余部分一个字不动：
%
% \begin{longtblr}[caption={三个旧类名的改法}, label={tab:oldclass}]
%   {colspec={l X[2,l]}, width=\linewidth}
% \toprule
%   旧写法 & 改成 \\ \midrule
%   \cs{documentclass}\marg{hithesisbook} &
%     \cs{documentclass}\oarg{stage=final,\ldots}\marg{hithesis} \\
%   \cs{documentclass}\marg{hithesisart} &
%     \cs{documentclass}\oarg{stage=proposal,\ldots}\marg{hithesis} \\
%   \cs{documentclass}\marg{hithesisartplus} &
%     \cs{documentclass}\oarg{stage=interim,degree-level=doctor,campus=shenzhen}\marg{hithesis} \\
% \bottomrule
% \end{longtblr}
%
%   写了旧名会收到一条说清改法的错误，不是 \LaTeX{} 自己那句看不出该改成什么的
%   \texttt{File~not~found}。
%
%   深圳博士中期原先自成一个文档类 \file{hithesisartplus.cls}，与
%   \file{hithesisart.cls} 九成代码逐字相同，现在是报告里
%   \texttt{campus=shenzhen,degree-level=doctor,stage=interim} 这一格，与其余十七格并列。
%   空壳会补齐这三个选项，所以不带选项的旧写法也照编。
%
%   三个空壳都在 v4（预计 2027 年发布）移除，届时不改的文档会直接编不过。
% \item \textbf{目录条目与 PDF 书签里的“结论”不再带撑开的空隙。}
%   规范要求“结论”“致谢”“目录”这类两字标题撑开排，原先撑开写在固定词汇的
%   取值里，同一个值又供目录条目与书签用，书签里读出来带着空隙。现在取值只留
%   文字，撑开由排标题的地方做。变的是目录条目、页眉与书签。
%
%   撑多少由 \option{term} 族的 \option{heading-spread} 定，默认一个全角字宽
%   （\texttt{1em}）——Word 版模板里那两个字之间是两个半角空格，宋体的半角空格
%   是半个字宽。原先写的是 \cs{ccwd}，那是“字宽加字距”，开着 \option{zijv}
%   时比一个字宽多 5.6\%，所以三个标题的字间距各窄了 1bp。博士后那一档规范
%   不撑开，取空值。
% \item \textbf{博士后封面的四个字段标签现在真的对齐了。}
%   “博士后姓名”“合作导师”“流动站名称”“专业名称”本该是同一列宽，原先靠手写
%   字间距凑。凑的时候字数按字宽算、间距按 \cs{ccwd} 算，而开了 \option{zijv}
%   之后 \cs{ccwd} 是字宽的 1.0558 倍，两条式子把这个差额放大成不同的量，
%   四个标签一个 99.94bp、一个 100.72bp，值那一列对不齐。现在改成分散对齐到
%   7 个字宽，四行都是 98.00bp，比原先窄 1.5bp。
% \item \textbf{“摘要”标题与“结论”“致谢”“目录”同一条路了。}
%   它原先也是把撑开写在取值里（\cs{cabstractcname}），同样漏进目录条目与书签。
%   现在一并走 \option{heading-spread}，字间距一个全角字宽，目录条目与书签取原文。
% \item \textbf{双语题注的 \cs{caption}\oarg{短标题} 与 \cs{label} 都能用了。}
%   \cs{caption}\oarg{短标题}\marg{中文题}\oarg{英文题} 里的短标题原先没进图清单，
%   而是被悄悄当成了 \cs{label}——\pkg{ccaption} 的 \cs{bicaption} 可选参数就是标签，
%   模板接错了位置。现在短标题进索引，与只写中文题时一致；\cs{label} 照普通
%   \hologo{LaTeX} 的写法放在 \cs{caption} 之后。
%
%   靠 \cs{caption}\oarg{名字} 当标签用的文档要改成 \cs{label}\marg{名字}，
%   否则那条引用会变成 \texttt{??}。
% \item \textbf{中文里的 \texttt{□}、\texttt{○}、\texttt{△} 排得出来了。}
%   \pkg{xeCJK} 按 Unicode 区块分中西文，“几何图形”那一块归了西文，所以中文正文里
%   写 \texttt{□}（规范原文用它占位，如“攻读□士学位”）会去找西文字体。
%   \pkg{TeX~Gyre~Termes} 里没有这个字，结果不是排成豆腐块，是整个字没了，
%   只在日志里留一句 \texttt{Missing character}。现在这六个符号归中文字体，正常排出。
% \item \textbf{\file{hithesis.sty} 现在只放示例自己的扩展。}
%   \cs{usepackage}\marg{hithesis} 仍可保留。它加载数学记号、\pkg{siunitx}、
%   \pkg{algorithm2e} 等示例功能；类排版必需的代码已经全部放回
%   \file{hithesis.cls}。不需要这些示例扩展时才删这一行。
% \item \textbf{\cs{bicaption} 不再需要。}
%   英文题注改成 \cs{caption} 的尾随可选参数，写
%   \cs{caption}\marg{中文题}\oarg{英文题}。\cs{bicaption} 的三种旧写法（\pkg{ccaption}
%   的五参数原形、本模板的三参数与四参数）都还能用，各发一条建议改写的警告。
% \item \textbf{子图从 \pkg{subfigure} 换成了 \pkg{subcaption}。}
%   \cs{subfigure}、\cs{subtable}、\cs{subref} 三个旧命令仍然能用，
%   \cs{ifsubcaphang} 等五个开关会自动转成对应的 \cs{captionsetup}，但都会发一条
%   建议改写的警告。
% \item \textbf{定理从 \pkg{ntheorem} 换成了 \pkg{thmtools}。}
%   \cs{theorembodyfont}、\cs{theoremheaderfont}、\cs{theoremsymbol}、
%   \cs{theoremseparator} 会转成 \cs{hitsetup} 里的对应键；\cs{theoremprework}、
%   \cs{theorempostwork}、\cs{theoremclass}、\cs{listtheorems}、\cs{thref} 没有
%   对应实现，会提示改用 \pkg{thmtools} 的写法。
% \item \textbf{类选项从 \pkg{kvoptions} 换成了 l3keys。}
%   排版类的选项写在 \cs{documentclass} 的可选参数里仍然生效，但会提示改用
%   \cs{hitsetup}，这条支持将在 v4 移除。
% \item \textbf{\env{publist} 改用 \pkg{enumitem} 实现。}
%   用法与版面都没变，仍然是 \cs{item} 逐条列。
% \item \textbf{\option{citetwo}\texttt{=comma} 现在真的生效了。}
%   这个选项从 2022 年加进来就一直是坏的（补丁打不上但吞掉了失败），修好之后
%   两个相邻连续序号会排成 \texttt{[1,2]} 而不是 \texttt{[1-2]}。原先误以为它
%   没作用的，升级后版面会变。
% \item \textbf{\cs{hitsetup}\marg{bib=\{adjacent-cite=…\}} 与
%   \cs{hitsetup}\marg{styles=\{Normal=…\}} 的格式项原先会报错}，现在修好了。
% \item \textbf{中文的斜体改成倾斜，不再是楷体。}
%   \cs{textit} 与 \cs{emph} 下的中文原先显示为楷体，现在是倾斜的正文字体。
%   要楷体请直接写 \cs{kaishu}。只影响写了 \option{font-zh} 的文档，不写这个选项
%   的中文字体仍由 \CTeX{} 设置，行为不变。
% \item \textbf{\option{font} 与 \option{cjk-font} 改名 \option{font-en} 与
%   \option{font-zh}。}
%   与模板其余选项的中英后缀一致。旧名仍然生效，只在日志里提示一次，将在 v4 移除。
% \item \textbf{\texttt{@masterthesis}（漏 \texttt{s} 的手误拼法）继续容错。}
%   \hologo{BibTeX} 的标准条目类型是 \texttt{@mastersthesis}。写漏 \texttt{s}
%   的拼法旧版一直认（深圳版曾经不认，v3.2a 开发期补齐过一轮），换官方底座后
%   官方原版不认，补丁里把这条容错带了过去，那条文献不会不声不响地消失。
% \item \textbf{报告的脚注改成了与终稿一致的样式。}
%   编号由上标阿拉伯数字改成带圈数字，正文改成宋体小五，悬挂缩进 1.5 字符，分割线
%   宽 0.3\cs{textwidth}。我工规范对各个类的脚注要求本来就是一样的，原先只有
%   终稿照做。
% \end{itemize}
%
% \subsection{新增的}
%
% \begin{itemize}
% \item \option{font-en} 与 \option{font-zh} 两个选项，西文与中文字体可以分别指定，
%   见第 \ref{sec:chinese-fonts} 节。
% \item \option{math-font} 与 \option{math-style} 两个选项，数学字体走
%   \pkg{unicode-math}，见第 \ref{sec:math-fonts} 节。
% \item \env{eqdenote} 环境，排公式下面的“式中……”注释表，见第
%   \ref{sec:eqdenote} 节。
% \item \option{layout/hyphenation} 选项，控制西文自动断字，取 \texttt{automatic} 或
%   \texttt{none}，默认 \texttt{none}。正文、题注、
%   中英文摘要一律跟随这一个开关。原文里本来就有的连字符（\texttt{Tree-structured}
%   这种）该断还断，不受它管。
% \item \env{publist} 列表在终稿与报告中使用同一套接口。
% \end{itemize}
%
% \chapter{使用说明}
% \label{sec:usage}
% 本手册假定用户已经能处理一般的 \hologo{LaTeX}\ 文档，并对 \hologo{BibTeX}\ 有一定了解。如果
% 从来没有接触过 \hologo{TeX}\ 和 \hologo{LaTeX}，建议先学习相关的基础知识。
%
% \section{提问}
% \label{sec:howtoask}
% 缺陷与可复现的问题请发到
% \href{https://github.com/hithesis/hithesis/issues}{GitHub Issues}。交流也可使用
% QQ 群 259959600、851792460、704864357。
%
% 提问时附上最小示例、完整日志、\hologo{TeX} Live 版本和编译引擎。若问题来自
% 院系通知或评审意见，请一并给出原始文件或可复核的截图；只转述“老师说不对”
% 无法判断该改哪一套版式。
%
% \section{示例文件}
% \label{sec:userguide}
% 新用户应从 \file{examples/demo/} 开始。\file{final.tex} 是中英文学位论文主文件，
% \file{report.tex} 是开题与中期报告主文件，\file{info.tex} 保存两者共用的元信息。
% 示例已经按前置、正文和后置材料分好目录，替换内容即可；不熟悉的接口再回本手册查。
%
% \lstinputlisting[style=lstStyleLaTeX]{examples/demo/final.tex}
%
% \section{示例宏包 \file{hithesis.sty}}
% \label{sec:example-package}
% v3.1e 的示例目录里有个 \file{hithesis.sty}，开头写着「此文件声明不在规范中要求的
% 格式所使用的宏包。（所以，格式基本上是自由发挥的。）」。里面装的是示例文档自己
% 要用的东西，跟排版规范没有关系，用户复制一份改成自己的。
%
% 分界线是“类自己调不调”。\pkg{siunitx}、\pkg{bm}、\pkg{mathrsfs}、\pkg{rotating}、
% \pkg{listings}、\pkg{algorithm2e} 这六个宏包，类里一句都没用到，只是替文档预先
% 装上并配好，全部归这个文件。类里留下的是自己要用的那些，比如封面用 \pkg{xcolor}
% 与 \pkg{tikz}，下划线用 \pkg{xeCJKfntef}。
%
% 里面装的东西：
% \begin{itemize}
% \item \pkg{siunitx} 与两句 \cs{sisetup}。规范对数字书写的规定是 4 位及以上每 3 位
%   空 1/4 汉字，这两句让 \cs{num} 照着排。除这个空格以外任何空格都是错的，
%   包括 \verb|\:| \verb|\;| \verb|\ |。
% \item 九个数学记号速记 \cs{theVector}、\cs{theMatrix}、\cs{theSet}、
%   \cs{theDirected}、\cs{theUndirected}、\cs{theNetwork}、\cs{theNode}、
%   \cs{theDirectedEdge}、\cs{theUndirectedEdge}。想换全文表示网络的符号样式，
%   改这一行就行，不用去正文里挨个找。\cs{theNetwork} 走 \pkg{rsfs} 花体，那套
%   字形只有 A 到 Z，小写字母和数字写进去会被悄悄丢掉。
% \item 整段 \pkg{algorithm2e} 配置：加载顺序、中英双语算法标题 \cs{AlgoBiCaption}、
%   算法列表 \texttt{algoen}、\cs{foocaption}。开题与中期报告是 \cls{article} 底子，
%   没有 \cs{chapter} 计数器，这段整个跳过，那边要排算法自己
%   \cs{usepackage}\marg{algorithm2e}。
% \item 六条示例缩略语词条，填在 \option{abbr} 族里。v3.1e 是每个中文术语手搓一个
%   宏（\cs{gtssbp}、\cs{gscna}、\cs{gscnas}）把 \cs{gls} 与两句 \cs{sindex} 打包，
%   现在正文一律写 \cs{abbr} 与 \cs{abbrs}，索引由词条自己带。
% \item 画染色体示意图的 \cs{gseg}、\cs{rcell}、\cs{tcellone} 等九个宏和配套的
%   \cs{tikzset} 样式。画的是固定内容，\cs{rcell} 直接画 A~C~T~C 四个碱基，
%   留在这里当画图的例子看。
% \item \cs{cs} 与 \cs{cmd}，在正文里排控制序列的名字。
% \end{itemize}
%
% \pkg{bm} 与 \pkg{mathrsfs} 看着跟类重复，其实不是。类只提供 \cs{boldsymbol}，
% 那不是 \cs{bm} 的全等替代；\cs{mathscr} 类里虽然有，但不装 \pkg{mathrsfs} 时它跟
% \cs{mathcal} 出自同一套字形，\cs{theNetwork} 与 \cs{theSet} 会长得一模一样，
% v3.1e 特意用 \pkg{rsfs} 就是为了把这两个记号分开。原文件里的 \pkg{xeCJKfntef}、
% \pkg{xcolor}、\pkg{tikz} 不在这里，那三个类自己要用，留在类里。
%
% 这些都不是规范要求的，改坏了也只影响自己的文档。改法是把
% \file{examples/demo/hithesis.sty} 拷进自己的项目目录再改，当前目录的那份优先于
% \hologo{TeX} 发行版里装的。
%
% 开题与中期报告的示例也加载它，主要为了共用正文中的数学记号与数字格式。
%
% \section{交付物：\option{stage}}
% \label{sec:class}
%
% 只有一个文档类 \file{hithesis.cls}。学位论文终稿、开题报告、中期报告都是它，
% 由 \option{stage} 选：
%
% \DescribeMacro{stage}
% \begin{latex}
% \documentclass[stage=final,    degree-level=doctor, campus=harbin]{hithesis}
% \documentclass[stage=proposal, degree-level=doctor, campus=harbin]{hithesis}
% \documentclass[stage=interim,  degree-level=doctor, campus=harbin]{hithesis}
% \end{latex}
%
% 三个取值：\texttt{final}（默认，学位论文终稿）、\texttt{proposal}（开题报告）、
% \texttt{interim}（中期报告）。不写就是 \texttt{final}。取值写错时日志里说一句，
% 按 \texttt{final} 排——这一步不能硬失败，类还没装完就报错的话，后面整份文档
% 都没有类可用。
%
% \option{degree-level}\texttt{=postdoc} 是个例外。博士后只写研究报告，没有开题与中期，
% 所以只要出现这个取值就按 \texttt{final} 排，\option{stage} 写什么都不看，不写也行。
% 显式写了别的 \option{stage} 时日志里说一句。
%
% \emph{换一个 \option{stage} 重编，正文里能用的东西不变。} 画图（\pkg{tikz}）、
% 颜色、三线表与长表（\pkg{tabularray} 的 \env{booktabs}、\env{longtblr}）、
% 索引项 \cs{hitindex}、缩略语 \cs{abbr} 与 \cs{abbrs}、双语题注的
% \option{caption-en}，这些两档都有。v3.2a 之前它们只在终稿那一档加载，同一份
% 正文改个 \option{stage} 会报一串未定义控制序列。
%
% \emph{整页级的构件只有终稿有，报告里静默不排。} \env{conclusions}、
% \env{acknowledgements}、\env{resume}、\env{publication}、\env{pubgroup}、
% \env{correspondingaddr}、\env{doctorpublication}、\env{ceindex} 八个环境，
% 学校的开题与中期报告规范里没有对应的一页。在报告那一档写了，整段吞掉，
% 不排也不报错——一份正文通吃三档，写了这几段是常态，不是手滑，不该逼你
% 把十来行调用注释来注释去。
%
% \cs{hitsetup} 顶层的 \option{cover}、\option{titlepage}、\option{defense}
% 三个键族不在此列：它们是\emph{设置}不是正文，在报告里写了会收到一条提示，
% 因为那多半真是写错了地方。
%
% 开题、中期和学位论文可以放在同一个目录里，共用元信息、参考文献库和
% 图片。三份主文件只需各自声明对应的 \option{stage}。
%
% \option{stage} 只能写在 \cs{documentclass} 上，挪不进导言区：\cs{LoadClass}
% 得在类文件里执行完，而 \cs{hitsetup} 是导言区的事，那时候基础类早装上了。
% 所以做不到“一个主文件靠一行开关切三种交付物”，能做到的是一个目录里三个
% 主文件，每个只有第一行不同。
%
% \section{论文选项}
% \label{sec:option}
%
% 类选项写在 \cs{documentclass} 的方括号里，示例见上节。这几项决定基础版面，
% 必须在加载基础类之前定下来，所以只能写在 \cs{documentclass} 里；
% 其余配置一律走 \cs{hitsetup}。
%
% \begin{longtblr}[caption={类选项}, label={tab:classopt}]
%   {colspec={l X[1,l] X[1,l]}, width=\linewidth}
% \toprule
% 选项 & 取值 & 管什么 \\ \midrule
% \option{degree-level}   & \texttt{bachelor} / \texttt{master} / \texttt{doctor} /
%                    \texttt{postdoc} & 学位，本科加星号是毕业设计 \\
% \option{degree-type} & \texttt{academic} / \texttt{professional} &
%                    学术学位还是专业学位，默认 \texttt{academic} \\
% \option{campus}  & \texttt{harbin} / \texttt{shenzhen} / \texttt{weihai} & 校区 \\
% \option{category} & \texttt{stem} / \texttt{hass} & 理工类还是人文社科类，默认
%                    \texttt{stem}，见下 \\
% \option{stage}    & \texttt{final} / \texttt{proposal} / \texttt{interim} &
%                    交付物种类，见 \ref{sec:class} 节 \\
% \option{lang}    & \texttt{zh} / \texttt{en} & 正文语言，默认 \texttt{zh} \\
% \option{fontset} & 见 \CTeX\ 手册 & 字体集 \\
% \bottomrule
% \end{longtblr}
%
% \option{lang} 只管摘要之后的正文。封面、扉页、原创性声明这些前置页面无论哪种
% 语言都照中文格式排，英文封面本来就是中文论文也有的那一页。
%
% \textbf{本科的毕业论文与毕业设计。} 本科封面上这是一对勾选框，写在
% \option{degree-level} 的取值里：
% \begin{latex}
% \documentclass[degree-level=bachelor,  campus=harbin]{hithesis}  % 毕业论文
% \documentclass[degree-level=bachelor*, campus=harbin]{hithesis}  % 毕业设计
% \end{latex}
% 星号只对本科有意义，别的学位写了不报错，但那对勾选框只出现在本科封面上。
%
% v3.2a 之前这件事由单独的 \option{design} 选项管。旧选项照旧能用，写了会在
% 日志里提示一次，兼容期到 v4（预计~2027~年发布）。
%
% \DescribeOption{category}
% 学校的写作指南与书写范例各有理工类、人文社科类两套，\option{category} 取
% \texttt{stem}（默认）或 \texttt{hass} 选哪一套。两份指南的打印要求逐条一样，
% 差在正文层次的编号：人文是“第一章”“一、”“（一）”“1.”，章序用一、二、三；
% 公式、图、表仍按阿拉伯章号编（指南原话：如第一章第一个公式的序号为（1-1））。
% 范例上又各自多出几处实现差别，人文这一档照两份人文范例：正文与目录固定值
% 23 磅、不贴网格；文献表半角“[1]”加一个空格、悬挂 1.71 字；成果页组名走节标题
% 那一档（黑体小三、段前段后 9.75 磅、编号“一、”）；英文目录里“Abstract(In
% Chinese)”“Acknowledgments”“letter of authorization”几处字样。
% \begin{latex}
% \documentclass[degree-level=doctor,campus=harbin,category=hass]{hithesis}
% \end{latex}
% \cs{hitsetup} 的条件也认它：\cs{hitsetup}\oarg{category=hass}\marg{...}。
% 章节命令的写法不变，编号的样子由这一档定。
%
% \DescribeOption{degree-type}
% \option{degree-level} 管学士硕士博士，学术还是专业是另一个维度，所以单开一个
% \option{degree-type}，取 \texttt{academic}（默认）或 \texttt{professional}。
%
% 封面上那行学生类别按它选文本：学术学位排 \option{student-type-academic}，
% 专业学位排 \option{student-type-professional}。这两条是 \option{term} 族的常量，
% 中英各一份，默认分别是“学术学位论文”“专业学位论文”与 \texttt{Academic Degree}
%“\texttt{Professional Degree}”。要改就改这两条：
% \begin{latex}
% \documentclass[degree-level=master,campus=harbin,degree-type=professional]{hithesis}
% ...
% \hitsetup{term={zh={student-type-professional={工程硕士}}}}
% \end{latex}
% 外面那对括号由模板加，值里不要自己带。
%
% \DescribeOption{degree-type}
% \option{degree-type} 这个元信息字段仍然认，作用是把上面选出来的文本整个换掉，
% 留给两条都不合适的情形，比如同等学力人员。不写就是空的，按 \option{degree-type}
% 走常量。
%
% v3.2a 的开发过程中这两条默认文本一度是往 \option{degree-type} 里写的，那是个错：
% 模板给的默认值与用户自己写的值挤在同一个槽里，分不出这个字段是空着还是被默认值
% 填过。示例文件里正好有一句 \option{degree-type}\texttt{=\{学术学位论文\}}，于是
% 写了 \option{degree-type}\texttt{=professional} 也看不出变化。默认值与覆盖入口
% 分开之后不会再有这个问题。
%
% 两档都认这个选项。报告目前没有按它分岔的版式，声明它是为了一份配置供两档
% 共用时条件写得出 \option{degree-type}\texttt{=professional}，以及以后规范真
% 分岔时有地方挂。
%
% 封面上这行什么时候排：新版面照 iota，研究生都排、本科不排；旧版面照 v3.1e，
% 非全日制、深圳、硕士这几种情形才排。写了 \option{degree-type} 不会让它多排一行，
% 要强制排或不排，写 \option{cover}\texttt{=\{degree-type=true\}} 或 \texttt{false}。
%
% \section{配置项的组织}
% \label{sec:keyfamily}
% \cs{hitsetup} 的配置项按职责分族。终稿收十三个族，报告收其中的八个：
%
% \begin{longtblr}[caption={\cs{hitsetup} 的配置族}, label={tab:family}]
%   {colspec={l X[1,l]}, width=\linewidth}
% \toprule
% 族 & 管什么 \\ \midrule
% \option{info}  & 封面与声明用到的元信息，如标题、作者、导师 \\
% \option{layout} & 版面，如纸张、页边距、文档网格、右开页 \\
% \option{styles} & 各处的字体字号行距间距 \\
% \option{toc}   & 目录，如排到第几级、目录里的西文字体 \\
% \option{bib}   & 参考文献，如后端（bibtex/biber）、样式名、著者大小写 \\
% \option{term} & 词条（封面标签、章节名这类固定词汇），见第~\ref{sec:const}~节 \\
% \option{math}  & 公式编号的括号形式 \\
% \option{struct} & 排哪些部件、各用哪个文件，见第~\ref{sec:structure}~节 \\
% \option{thm}   & 定理环境的字体与间距，见第~\ref{sec:math}~节 \\
% \option{cover} & 封面一上几处按学位分岔的取舍 \\
% \option{titlepage} & 内封上几处按学位分岔的取舍 \\
% \option{defense} & 答辩决议上的委员名单 \\
% \option{abbr}  & 缩略语词条，见第~\ref{sec:abbr}~节 \\
% \bottomrule
% \end{longtblr}
%
% 报告没有封面一、内封、答辩决议、缩略语表这几样，所以那边只收
% \option{info}、\option{layout}、\option{styles}、\option{toc}、\option{bib}、
% \option{term}、\option{struct}、\option{thm} 八个。
%
% 写法是嵌套的键值表：
% \begin{latex}
% \hitsetup{
%   info = {
%     zh = { title = {论文标题}, author = {张三} },
%   },
%   toc = { english = true },
%   styles = { Heading 1 = { font = { font-style = regular } } },
% }
% \end{latex}
%
% 族名之下的键用连字符分词，如 \option{degree-level}、\option{adjacent-cite}。
%
% \option{toc} 族一共六个键：
%
% \begin{longtblr}[caption={\option{toc} 族的键}, label={tab:toc}]
%   {colspec={l X[1,l] X[2,l]}, width=\linewidth}
% \toprule
% 键 & 取值 & 管什么 \\ \midrule
% \option{depth} & 整数，默认 \texttt{3} &
%   目录排到第几级。规范里用得到的是 \texttt{3} 与 \texttt{4}，
%   文科生要四级 \\
% \option{english} & 三态，默认 \texttt{auto} &
%   中文目录之后再排不排一份英文目录。\texttt{auto} 是博士排、其余不排；
%   英文版（\option{lang}\texttt{=en}）只有这一份目录，\texttt{auto} 一律排 \\
% \option{latin-natural} & 三态，默认 \texttt{auto} &
%   纯西文条目按不按西文字体自己的行高。\texttt{auto} 为真，是 Word 的原样：
%   只有西文的一条比汉字行矮一点；关掉就一律按汉字行高 \\
% \option{latin-sans} & 布尔，默认 \texttt{false} &
%   目录里的西文用不用无衬线体 \\
% \option{number-align} & \texttt{natural}\,/\,\texttt{left}\,/\,\texttt{right}，默认 \texttt{natural} &
%   条目编号怎么摆。\texttt{natural} 排自然宽、编号后空一个字（范例的排法）；
%   \texttt{left} 是量宽取齐的旧行为；\texttt{right} 编号右贴列缘 \\
% \option{lists} & 布尔，默认 \texttt{false} &
%   图表清单要不要进目录。规范没有这一项，范例的目录里也没有这两条，
%   默认不进；写 \texttt{true} 中英两份目录一起进 \\
% \bottomrule
% \end{longtblr}
%
% 章节标题里的西文归 \option{styles}~=~\{~\option{heading}~=~\{~\option{font}~=~\{~\option{latin-text-font}~=~\texttt{Arial}~\}~\}~\}，
% 与 \option{toc}~=~\{~\option{latin-sans}~\} 是一对。
%
% \textbf{英文目录里的条目从哪来。}四级章节命令的末尾可选参数是英文名：
% \cs{chapter}\marg{中文标题}\oarg{English Title}。给了就照它排；没给不再
% 整条消失，退回排印用的标题，中文版再在后面缀一句红字
% \texttt{(Error! English heading not defined.)}，句式照 Word 的“错误！未定义
% 书签。”，只提示不拦编译，补上英文名就消失。英文版
% （\option{lang}\texttt{=en}）的标题本来就是英文，退回就是它，不缀提示。
% 那句字样是词条 \texttt{table-of-contents-undefined-en}，\option{term}
% 改得动。
%
% \option{depth} 与 \option{latin-sans} 两档都认。\option{depth} 数的是级数，
% 不是 \cs{tocdepth}：学位论文最高一级是章，报告是节，两边换算不一样，
% 写级数就不用管这个差别。\option{english} 只有终稿有，
% 报告一份文档只排一份目录，语言跟着类选项 \option{lang} 走，
% 在报告里写它会收到一条提示。
%
% \option{english} 说的是排几份目录，跟 \option{struct} 族的
% \option{table-of-contents} 不是一回事：\option{table-of-contents} 决定这一份文档排不排目录，
% \option{english} 决定排出来的目录里有没有英文那一份。
%
% 报告还认一个 \option{toc}~=~\texttt{true}\,/\,\texttt{false}，它是
% \option{struct}~=~\{~\option{frontmatter}~=~\{~\option{table-of-contents}~\}~\} 的别名，
% 写哪个都行，两个不会各排一遍。学位论文没有这个简写，写 \option{table-of-contents}。
%
% 报告不写就有目录。其余部件都是没声明就不排，这一个相反：\option{toc} 是
% v3.1e 的类选项，默认开着，所以报告的 \option{table-of-contents} 默认为 \texttt{true}。
% 深圳博士中期是唯一的例外，那一格默认不排目录，显式写
% \option{toc}~=~\texttt{true} 才排。
%
% \subsection{一份配置供多种交付物}
% \label{sec:setup-condition}
% 一个目录里放学位论文、开题、中期三份主文件时，配置通常写在一份
% \file{hitsetup.tex} 里，由三份主文件各自 \cs{input}。可两档接受的键不完全一样：
% 声明页、授权书、致谢这些词只有终稿有，页眉那两条只有报告有。不加区分地
% 写在一起，各自会在日志里报几条“不认识的键”，跟真的敲错键名拿到同一条提示。
%
% 给 \cs{hitsetup} 加个可选参数，写条件：
%
% \begin{latex}
% \hitsetup[stage=final]{ term = { declarations-heading-zh = {使用授权书} } }
% \hitsetup[class=report]{ term = { header-report = {哈尔滨工业大学开题报告} } }
% \hitsetup{ info = { zh = { author = {张三} } } }   % 不带条件，都生效
% \end{latex}
%
% 条件不成立时整块跳过，连键都不去设，所以不会有任何提示。
%
% 条件的写法与取值：
%
% \begin{longtblr}[caption={\cs{hitsetup} 条件的写法}, label={tab:cond}]
%   {colspec={l X[l]}, width=\linewidth}
% \toprule
% 写法 & 意思 \\ \midrule
% \texttt{[stage=final]} & 终稿 \\
% \texttt{[class=report]} & 开题中期报告 \\
% \texttt{[stage=proposal]}、\texttt{[stage=interim]} & 报告里的开题、中期 \\
% \texttt{[campus=shenzhen]}、\texttt{[degree-level=doctor]} & 校区、学位 \\
% \texttt{[campus=harbin,degree-level=doctor]} & 逗号是“且” \\
% \texttt{[campus=weihai\textbar campus=shenzhen]} & 竖线是“或” \\
% \texttt{[campus=shenzhen;degree-level=doctor]} & 分号分组，组之间是“或” \\
% \bottomrule
% \end{longtblr}
%
% 等号左边只是给人看的，真正匹配的是右边那个取值对应的开关，所以
% \texttt{[foo=shenzhen]} 与 \texttt{[campus=shenzhen]} 效果一样。右边写了不存在的
% 取值会报一条“不认识的取值”，那一组当作不成立。
%
% \subsection{三态取值}
% \label{sec:tristate}
% 有些配置项的合适取值随学位而变，这类配置项取三种值：
%
% \begin{longtblr}[caption={三态取值}, label={tab:tristate}]
%   {colspec={l X[1,l]}, width=\linewidth}
% \toprule
% 取值 & 含义 \\ \midrule
% \texttt{auto} & 由模板按学位和校区决定，默认值 \\
% \texttt{true} / \texttt{false} & 一律开或一律关 \\
% 分学位映射 & 逐个学位指定，没列到的学位仍按 \texttt{auto} \\
% \bottomrule
% \end{longtblr}
%
% \begin{latex}
% \hitsetup{
%   toc   = { english = auto },                         % 模板决定
%   layout = { frontmatter = { page-number-line = auto } },
% }
% \end{latex}
%
% 常见的三态键有 \option{toc/english}，以及 \option{layout} 族中按部件设置的
% \option{openright} 与 \option{page-number-line}。
%
% \subsection{文档骨架}
% \label{sec:structure}
% \DescribeMacro{\hitfrontmatter}
% \DescribeMacro{\hitbackmatter}
% 前置与后置部分的排列顺序是规范定死的：附录在参考文献后，决议在成果后，
% 授权在索引后。用户没有自由度，却要自己在主文件里抄对二十来行 \cs{include}，
% 抄错了不报错，只是排错。
%
% 所以可以只说“哪个部件用哪个文件”，顺序交给模板：
%
% \emph{开关型部件的取值有三档。} \texttt{true} 排、\texttt{false} 不排，
% \emph{不写或写 \texttt{auto}} 是交给模板按类与选项定。三个索引的默认条件是
% “这本是不是英文版”——规范只对英文版要求排图索引与表索引；公式索引默认不排，
% 要就显式写 \texttt{true}。
%
% 名字取全称：\option{table-of-contents}、\option{list-of-figures}、
% \option{list-of-tables}、\option{list-of-equations}，读起来就是那一页的标题。
% \option{index} 是书末的主题索引，与这三个不是一类，名字照旧。
% \begin{latex}
% \hitsetup[stage=final,lang=zh]{ struct = {
%   frontmatter = {
%     cover      = {final/frontmatter/cover},
%     abstract   = {final/frontmatter/abstract},
%     list-of-symbols = {final/frontmatter/symbols},
%     list-of-abbreviations = true,
%     table-of-contents = true,
%     list-of-figures   = auto,   % 不写也是 auto
%     list-of-tables    = auto,
%   },
%   mainmatter = {final/mainmatter/ch1, final/mainmatter/ch2,
%                 final/mainmatter/conclusion},
%   backmatter = {
%     bibliography     = {final/backmatter/reference},
%     appendix         = {final/backmatter/appA},
%     publications     = {final/backmatter/publications},
%     defense          = auto,
%     index            = {final/backmatter/ceindex},
%     declarations     = true,
%     acknowledgements = {final/backmatter/acknowledgements},
%     resume           = {final/backmatter/resume},
%   },
% } }
% \begin{document}
% \end{document}
% \end{latex}
%
% \DescribeMacro{\hitmainmatter}
% \textbf{三个骨架宏一句都不用写。} \cs{hitfrontmatter}、\cs{hitmainmatter}、
% \cs{hitbackmatter} 没写的话，模板在文末按前置、正文、后置的次序补齐。
% 想在两段之间插点别的东西时再显式写出来，写了的就在写的位置排。
%
% 混着来要当心：自己写了正文却不写 \cs{hitmainmatter}，补出来的正文会落在
% 自己写的内容后面。
%
% \option{mainmatter} 取一串文件名，逗号分隔，\textbf{按写的顺序}逐个
% \cs{include}——正文各章的先后是作者的事，不是规范的事。结论是正文的最后一章，
% 写进这一串就行。逗号后面的空格不必省。
%
% \DescribeOption{frontmatter}
% \DescribeOption{backmatter}
% 部件写进 \option{frontmatter} 或 \option{backmatter} 分组，\textbf{分组决定这个
% 部件归哪一段，段内部的顺序仍由规范定}。顺序不交给用户，是因为这个特性本来就是
% 为了消灭“抄错顺序”这类错误。
%
% 能挪的部件其实只有致谢：中文论文排在后置，英文论文习惯排在前置。
% 封面与摘要压根不是“排在哪”，它们要赶在封面排版之前读进来；参考文献、附录、
% 索引、授权几样规范写死在后置，彼此还有先后约束。
%
% \begin{longtblr}[caption={\option{struct} 族的部件}, label={tab:struct}]
%   {colspec={l l l}}
% \toprule
% 键 & 取值 & 归哪一段 \\ \midrule
% \option{cover}            & 文件名 & 前置，只读不排 \\
% \option{abstract}         & 文件名 & 前置，只读不排 \\
% \option{list-of-symbols}  & 文件名 & 前置，目录之前 \\
% \option{list-of-abbreviations} & 三值 & 前置，符号表之后 \\
% \option{list-of-symbols-and-abbreviations} & 文件名 & 前置，两者排一页 \\
% \option{table-of-contents} & \texttt{auto} / \texttt{true} / \texttt{false} & 前置，最后。报告里默认开 \\
% \option{mainmatter}       & 文件名表，按写的顺序 & 正文 \\
% \option{bibliography}     & \texttt{.bib} 文件名 & 后置，第一项 \\
% \option{appendix}         & 文件名 & 后置，参考文献之后 \\
% \option{publications}     & 文件名 & 后置，附录之后 \\
% \option{defense}          & \texttt{auto} / \texttt{true} / \texttt{false} / 文件名 & 后置，成果之后 \\
% \option{defense*}         & 文件名 & 同上，贴 PDF 并排页码 \\
% \option{index}            & 文件名 & 后置，决议之后 \\
% \option{declarations}     & \texttt{true} / \texttt{false} / 文件名 & 后置，索引之后 \\
% \option{declarations*}    & 文件名 & 同上，贴 PDF 并排页码 \\
% \option{acknowledgements} & 文件名 & 后置或前置，可挪 \\
% \option{resume}           & 文件名 & 后置，最后 \\
% \bottomrule
% \end{longtblr}
%
% \option{cover} 与 \option{abstract} 那两个文件只是被读进来：封面那个设元信息，
% 摘要那个把两段正文收起来，真正排版由模板接着封面一起做。所以它们必须赶在
% 封面之前读到，写在分组里的哪个位置都不影响排出来的次序。
%
% 取文件名的那几项走 \cs{include}，只有 \option{cover}、\option{abstract}、
% \option{list-of-symbols} 与 \option{appendix} 走 \cs{input}，附录还会套上
% \env{appendix} 环境。
%
% 写错了会说三种话：
% \begin{itemize}
% \item 部件放进它不能去的那一段，报错并说清它能写在哪；
% \item 同一次 \cs{hitsetup} 里把一个部件写进两个段，警告，按后写的排；
% \item 分组里出现不认识的键，警告并忽略。
% \end{itemize}
%
% 报告也认这几个宏，这样一个目录里几份交付物的正文部分长得一样。报告没有
% 前后文之分，\cs{hitfrontmatter} 只排封面，\cs{hitbackmatter} 只排参考文献，
% \option{mainmatter} 走 \cs{input} 而不是 \cs{include}——报告正文紧接目录排，
% 多一个分页就多一张空白页。
%
% 原先逐个 \cs{include} 的写法照旧有效，两种可以混着用。
%
% \subsection{一份配置管几种交付物}
% \label{sec:one-main}
% \option{stage} 必须写在 \cs{documentclass} 上，因此每种交付物仍各有一个很薄的
% 主文件。元信息与结构配置可以放进同一份文件，由 \cs{hitsetup} 的条件参数分流：
% \begin{latex}
% \hitsetup[stage=final,lang=zh]{ struct = { ...完整一份... } }
% \hitsetup[stage=final,lang=en]{ struct = { ...完整一份... } }
% \hitsetup[stage=proposal]{ struct = { ...完整一份... } }
% \hitsetup[stage=interim]{ struct = { ...完整一份... } }
% \end{latex}
%
% \textbf{条件决定“做不做”，不决定“谁优先”。} \cs{hitsetup} 一贯是后写的覆盖
% 先写的，这一条不因为带了条件而改变。上面四组两两互斥，任何一份文档只命中一组，
% 所以先后无关。
%
% 写成“一组打底、另一组只改几项”也可以，但那样两组会同时命中，
% \textbf{顺序就有讲究了}：打底的那组必须写在前面，写反了它会把改过的项盖回去，
% 而且不报错。
%
% \subsection{按交付物加载宏包}
% \label{sec:usepackage}
% \DescribeMacro{\hitusepackage}
% 一份主文件只有一个导言区，\cs{usepackage} 对哪种交付物都生效。只有某一种要的
% 宏包写成：
% \begin{latex}
% \hitusepackage[stage=proposal]{lipsum, siunitx={per-mode=symbol, detect-all}}
% \end{latex}
% 条件与 \cs{hitsetup} 的那个可选参数完全一样，逗号是且、竖线是或、分号分组。
% 带选项的写成\meta{包名}\texttt{=}\marg{选项}：选项用花括号不用方括号，
% 外层是逗号表，方括号里的逗号会被外层拆开，花括号护得住。
%
% 按写的顺序加载，不排序也不去重——宏包之间的先后常常有讲究。
%
% 只能写在导言区，也就是 \cs{begin}\marg{document} 之前。
%
% \subsubsection{文档类已经装好的和没装的}
% 文档类只装它自己排版要用的宏包。示例目录下的 \file{hithesis.sty} 另装一批
% 常用的：\pkg{siunitx}、\pkg{bm}、\pkg{mathrsfs}、\pkg{rotating}、\pkg{listings}、
% \pkg{algorithm2e}、\pkg{multirow}、\pkg{placeins}。
% 导言区写 \cs{usepackage}\marg{hithesis} 就都有了，示例的 \file{final.tex} 里
% 已经写好。
%
% \pkg{multirow} 与 \pkg{placeins} 在 v3.2a 从文档类挪到示例宏包。老文档若用了
% \cs{multirow} 或 \cs{FloatBarrier}，应保留 \cs{usepackage}\marg{hithesis}，
% 或者自行加载对应宏包。\pkg{xeCJKfntef} 仍由类按引擎加载。
%
% 别的宏包自己装。只给某一种交付物用的写 \cs{hitusepackage}，
% 每种交付物都要的写 \cs{usepackage}。
%
% \subsubsection{\texttt{auto} 的处理}
% \option{struct} 里的部件分两类。\option{table-of-contents}、\option{defense}、
% \option{declarations} 是开关，写 \texttt{auto} 的意思是“让模板按规矩决定”：
% 决议只有博士排，目录与授权书都排。其余部件要读一个文件，写 \texttt{auto}
% 没有意义，模板会说一句并跳过这一块，而不是去找 \file{auto.tex}。
%
% 开关型部件拿不准时可以写 \texttt{auto}。需要文件名的部件不要这么写；
% 模板会警告并跳过它。
%
% \subsection{答辩决议那一页}
% \label{sec:defense}
% 这一页排的是评阅人名单、答辩委员会名单和答辩委员会决议。默认只有博士排，
% 由 \option{struct} 里的 \option{defense} 控制：
%
% \begin{longtblr}[caption={\option{defense} 的取值}, label={tab:defense}]
%   {colspec={l l}}
% \toprule
% 取值 & 谁排 \\ \midrule
% \option{auto}  & 只有博士排 \\
% \option{true}  & 都排 \\
% \option{false} & 都不排 \\
% \bottomrule
% \end{longtblr}
%
% \DescribeOption{defense}
% 名单与决议正文写在 \option{defense} 族里，一个人一条记录：
% \begin{latex}
% \hitsetup{ defense = {
%   reviewer  = {name=张三, title=教授（博导）, affiliation=哈尔滨工业大学,
%                discipline=机械工程},
%   reviewer  = {name=李四, title=研究员, affiliation=中国科学院沈阳自动化研究所,
%                discipline=机械工程},
%   chair     = {name=王五, …},
%   member    = {name=赵六, …},
%   member    = {name=孙七, …},
%   secretary = {name=周八, …},
%   resolution = {答辩委员会听取了论文作者的报告……建议授予工学博士学位。},
% } }
% \end{latex}
%
% 五个键：\option{reviewer} 评阅人，\option{chair} 主席，\option{member} 委员，
% \option{secretary} 秘书，\option{resolution} 决议正文。一条记录里四个字段：
% \option{name}、\option{title}、\option{affiliation}、\option{discipline}，
% 分别对应表上的姓名、职称（是否博导）、工作单位、所在学科。
%
% \textbf{同一个键写多遍是累加，不是覆盖。} 评阅人与委员写几条就是几个人，
% 顺序就是写的顺序。某个字段不写只是那一格空着，不会把后面的顶错位。
% 键与键之间空一行分组也行。
%
% 一条都不填也能编，排出来是一张空表，行数与纸质表一样，打印出来手写或盖章。
% 填了人就按人数排，比纸质表的行数多也行。
%
% \textbf{行数会互相匀。} 纸质表四块的行数是评阅人 2、主席 1、委员 6、秘书 1。
% 某一块写超了，先看别的块有没有空着的行——委员只填四个就空两行，这两行匀给
% 超了的那块，总行数不变，决议区也就不会变矮。全借完还不够才加总行数。
%
% \textbf{“委员”这类竖排标签会跟着行数变。} 那一块只剩一行时两个字自己横排，
% 行数够了变成竖排。两个字的标签字距随格子高度撑开，最多撑到 1.8 倍行距，
% 与纸质表默认六行时的间距一致；三个字以上的标签贴着排、整摞居中，
% 也与纸质表一致。
%
% 决议正文那一框吃掉整页剩下的高度：上面两张表有几个人就多高，剩多少这一框就占
% 多少，整页的底边始终落在纸质表的同一位置。
%
% 决议正文写得长、一页放不下，就接着排到下一页，续页那一框四条边都有，
% 线宽与第一页最外侧那圈一样。名单填得多、把决议的地方挤到放不下两行字时，
% 决议整块挪到下一页，第一页由名单表撑满，底边照旧落在原处。
% 这时行高会比纸质表的略高，是保住底边对齐的代价。
%
% \DescribeOption{defense*}
% 已经有盖过章的扫描件就把文件名写进同一个键，不必另外调命令：
% \begin{latex}
% struct = { backmatter = {
%   defense  = auto,                        % 只有博士排，排版式表格
%   defense  = {figures/resolution.pdf},    % 贴这个 PDF，不排页码
%   defense* = {figures/resolution.pdf},    % 贴这个 PDF，排页码
% } }
% \end{latex}
% 取值是 \texttt{auto}、\texttt{true}、\texttt{false} 三个词之一就按上面的数据排，
% 别的都当文件名。\option{declarations} 与 \option{declarations*} 写法完全一样。这跟 \option{bibliography}、\option{appendix}
% 那些“给文件名就读那个文件”是同一个路数。
%
% \DescribeMacro{\resolution}
% \cs{resolution}\oarg{文件} 与 \cs{resolution*}\oarg{文件} 照旧能用，效果与上面两行
% 相同，但已经是旧写法：调用时日志里会提示改写成 \option{struct} 的 \option{defense}，
% 兼容期到 v4（预计~2027~年发布）。
%
% \textbf{从 v3.1e 升上来要注意：} 不给文件时原先贴的是
% \file{figures/empty-resolution.pdf} 那张扫描的空表图，现在改成排出来，
% 那张图也从示例目录删掉了。自己项目里还留着那张图不影响编译，只是用不上了。
% 另外这个键从 v3.2a 起叫 \option{defense}（\option{resolution} 现在是
% \option{defense} 族里那个决议正文的键）；写 \texttt{true} 原先只有博士生效，现在谁写谁排，
% 想保持原样就改成 \option{auto}。
%
% \subsection{旧写法}
% \label{sec:deprecated}
% v3.2a 之前所有配置项都写在同一层，且有一批名字带 \texttt{bs} 前缀或不带连字符。
% 旧写法仍然可用，编译时会在日志里提示新写法，兼容期一年，
% v4（预计~2027~年发布）移除。对照表：
%
% \begin{longtblr}[caption={旧配置项与新写法对照}, label={tab:renamed}]
%   {colspec={l X[1,l]}, width=\linewidth}
% \toprule
% 旧名 & 新写法 \\ \midrule
% \option{arialtoc} & \option{toc} = \{ \option{latin-sans} \} \\
% \option{tocfour} & \option{toc} = \{ \option{depth} = \texttt{4} \} \\
% \option{tocblank} & 已废弃，见第~\ref{sec:removed}~节 \\
% \option{engtoc} & \option{toc} = \{ \option{english} \} \\
% \option{arialtitle} & \option{styles} = \{ \option{heading} = \{ \option{font} = \{ \option{latin-text-font} = \texttt{Arial} \} \} \} \\
% \option{chapterbold} & \option{styles} = \{ \option{Heading 1} = \{ \option{font} = \{ \option{font-style} = \texttt{bold} \} \} \} \\
% \option{chapterhang} & \option{styles} = \{ \option{Heading 1} = \{ \option{paragraph} = \{ \option{indentation} = \{ \option{special} = \texttt{hanging} \} \} \} \} \\
% \option{capcenterlast} & \option{styles} = \{ \option{caption} = \{ \option{paragraph} = \{ \option{alignment} = \texttt{last-line-centered} \} \} \} \\
% \option{subcapcenterlast} & \option{styles} = \{ \option{caption-sub} = \{ \option{paragraph} = \{ \option{alignment} = \texttt{last-line-centered} \} \} \} \\
% \option{absupper} & \option{styles} = \{ \option{abstract-heading} = \{ \option{font} = \{ \option{all-caps} \} \} \} \\
% \option{etitleauto} & \option{titlepage} = \{ \option{title-en-xiaoer} \} \\
% \option{glue} & 旧版面专用，新版面的行距是刚性的 \\
% \option{openright} & \option{layout} = \{ \option{openright} \} \\
% \option{library} & \option{layout} = \{ \option{openright} = \texttt{false} \} \\
% \option{raggedbottom} & \option{layout} = \{ \option{bottom} = \texttt{ragged} 或 \texttt{flush} \} \\
% \option{zijv} & 旧版面专用，新版面按字符网格 \\
% \option{bsheadrule} & \option{layout} = \{ \option{header} = \{ \option{rule} \} \} \\
% \option{bsmainpagenumberline} & \option{layout} = \{ \option{mainmatter} = \{ \option{page-number-line} \} \} \\
% \option{bsfrontpagenumberline} & \option{layout} = \{ \option{frontmatter} = \{ \option{page-number-line} \} \} \\
% \option{citetwo} & \option{bib} = \{ \option{compress-min} = \texttt{2}｜\texttt{3} \} \\
% \option{splitbibitem} & \option{bib} = \{ \option{split-entry} \} \\
% \option{biblabelalign} & \option{bib} = \{ \option{numbering-align} \} \\
% \option{bibauthorupper} & \option{bib} = \{ \option{uppercase-family} \} \\
% \option{bibmaxauthor} & \option{bib} = \{ \option{max-bib-names} \} \\
% \bottomrule
% \end{longtblr}
%
% 元信息字段在 v3.2a 也改成了英文名。原先的 \texttt{c}\ldots{} 与 \texttt{e}\ldots{}
% 前缀取消，中英文两版改用 \texttt{-zh} 与 \texttt{-en} 后缀，与 \option{font-zh}、
% \option{font-en} 一致。旧名同样保留到 v4。
%
% \begin{longtblr}[caption={元信息字段的旧名与新名对照}, label={tab:renamed-info}]
%   {colspec={l X[1,l]}, width=\linewidth}
% \toprule
% 旧名 & 新写法 \\ \midrule
% \option{caffil} & \option{info} = \{ \option{affiliation-zh} \} \\
% \option{eaffil} & \option{info} = \{ \option{affiliation-en} \} \\
% \option{cassosupervisor} & \option{info} = \{ \option{associate-supervisor-zh} \} \\
% \option{eassosupervisor} & \option{info} = \{ \option{associate-supervisor-en} \} \\
% \option{cauthor} & \option{info} = \{ \option{author-zh} \} \\
% \option{eauthor} & \option{info} = \{ \option{author-en} \} \\
% \option{cclassid} & \option{info} = \{ \option{class-id} \} \\
% \option{intclassifiedindex} & \option{info} = \{ \option{classified-index-international} \} \\
% \option{natclassifiedindex} & \option{info} = \{ \option{classified-index-national} \} \\
% \option{ccosupervisor} & \option{info} = \{ \option{co-supervisor-zh} \} \\
% \option{ecosupervisor} & \option{info} = \{ \option{co-supervisor-en} \} \\
% \option{firstpagecdate} & \option{info} = \{ \option{cover-date-zh} \} \\
% \option{ctitlecover} & \option{info} = \{ \option{cover-title-zh} \} \\
% \option{cdate} & \option{info} = \{ \option{date-zh} \} \\
% \option{edate} & \option{info} = \{ \option{date-en} \} \\
% \option{cxueke} & \option{info} = \{ \option{discipline-zh} \} \\
% \option{exueke} & \option{info} = \{ \option{discipline-en} \} \\
% \option{cenrolldate} & \option{info} = \{ \option{enroll-date-zh} \} \\
% \option{ckeywords} & \option{info} = \{ \option{keywords-zh} \} \\
% \option{ekeywords} & \option{info} = \{ \option{keywords-en} \} \\
% \option{cnumber} & \option{info} = \{ \option{report-number} \} \\
% \option{csubmitdate} & \option{info} = \{ \option{report-submit-date-zh} \} \\
% \option{ctype} & \option{info} = \{ \option{report-type-zh} \} \\
% \option{cenddate} & \option{info} = \{ \option{research-end-date-zh} \} \\
% \option{cstartdate} & \option{info} = \{ \option{research-start-date-zh} \} \\
% \option{statesecrets} & \option{info} = \{ \option{secret-level} \} \\
% \option{csubject} & \option{info} = \{ \option{major-zh} \} \\
% \option{esubject} & \option{info} = \{ \option{major-en} \} \\
% \option{cpositionname} & \option{info} = \{ \option{station-zh} \} \\
% \option{cstudentid} & \option{info} = \{ \option{student-id} \} \\
% \option{cstudenttype} & \option{info} = \{ \option{degree-type-zh} \} \\
% \option{estudenttype} & \option{info} = \{ \option{degree-type-en} \} \\
% \option{csubtitle} & \option{info} = \{ \option{subtitle-zh} \} \\
% \option{esubtitle} & \option{info} = \{ \option{subtitle-en} \} \\
% \option{csupervisor} & \option{info} = \{ \option{supervisor-zh} \} \\
% \option{esupervisor} & \option{info} = \{ \option{supervisor-en} \} \\
% \cs{bicaption}\marg{中文题}\marg{英文题} & \cs{caption}\marg{中文题}\oarg{英文题} \\
% \cs{longbionenumcaption}（七参数） & \cs{caption}\marg{中文题}\oarg{英文题} \\
% \env{cabstract} & \env{abstract-zh} \\
% \env{eabstract} & \env{abstract-en} \\
% \option{ctitle} & \option{info} = \{ \option{title-zh} \} \\
% \option{ctitlecover} & \option{info} = \{ \option{title-zh} \} \\
% \option{csubtitle} & \option{title-zh} 的 \oarg{副标题} \\
% \option{esubtitle} & \option{title-en} 的 \oarg{副标题} \\
% \option{etitle} & \option{info} = \{ \option{title-en} \} \\
% \option{ctitleone} & \option{title-zh} 的第一个 \marg{分组} \\
% \option{ctitletwo} & \option{title-zh} 的第二个 \marg{分组} \\
% \option{cfinishdate} & \option{info} = \{ \option{work-finish-date-zh} \} \\
% \bottomrule
% \end{longtblr}
%
% 排版类的配置项此前也可以写在 \cs{documentclass} 的可选参数里，同样进入兼容期，
% 请改用 \cs{hitsetup}。\cs{documentclass} 里保留 \option{degree-level}、\option{campus}、
% \option{fontset}、\option{lang} 与 \option{stage}，它们决定版面，必须在加载
% 基础类之前就定下来。
%
% \subsection{废弃的写法}
% \label{sec:removed}
% 上一节那些是改了名，这一节是没有新名可改的。模板仍然认这些键，写了只在日志里
% 说一句，什么也不做。
%
% \begin{longtblr}[caption={废弃的配置项}, label={tab:removed}]
%   {colspec={l X[1,l]}, width=\linewidth}
% \toprule
% 废弃的 & 为什么 \\ \midrule
% \option{tocblank} & 目录里第一章之前的那行空白规范中没有要求 \\
% \bottomrule
% \end{longtblr}
%
% \section{西文字体}
% 西文字体由类选项 \option{font-en} 决定：
%
% \begin{longtblr}[caption={\option{font-en} 的取值}, label={tab:font-en}]
%   {colspec={l X[1,l]}, width=\linewidth}
% \toprule
% 取值 & 字体 \\ \midrule
% \texttt{auto} & 默认。见下 \\
% \texttt{times} & Times New Roman + Arial + Courier New \\
% \texttt{termes} & TeX Gyre Termes + Heros + Cursor \\
% \texttt{none} & 不设置，交给用户自己写 \\
% \bottomrule
% \end{longtblr}
%
% \texttt{auto} 时按 \option{fontset} 决定：写的是开源字体集
% （\texttt{fandol}、\texttt{ubuntu}）就用 TeX Gyre Termes（Times 的开源克隆），
% 其余一律优先 Times New Roman——学校要求的就是它。机器上找不到会自动退回
% TeX Gyre，不会报错，所以 Linux 上装没装 Times New Roman 都能编。
%
% 之所以给开源字体集配开源西文，是因为写 \texttt{fandol} 或 \texttt{ubuntu} 的人
% 多半图的是“到哪台机器都编出一样的东西”，配一个装没装看运气的字体不合这个意思。
% 要在开源字体集上也用 Times New Roman，显式写 \option{font-en=times} 即可。
%
% \section{中文字体}
% \label{sec:chinese-fonts}
% 正确配置中文字体是使用模板的第一步。模板调用 \CTeX\ 宏包，只提供基于
% \pkg{xeCJK} 包，使用 \hologo{XeLaTeX}\ 编译的方式。
% 关于如何使用字体命令、字号等等，属于模板格式范畴，在实现细节中讨论。
% 关于中文字体安装、配置的所有问题不在本模板讨论 范围。
%
% 默认情况下中文字体由 \CTeX\ 按 \option{fontset} 配置，模板不插手。写了类选项
% \option{font-zh} 就改由模板自己设置，此时 \option{fontset} 只用来决定
% \option{font-zh} 取 \texttt{auto} 时选哪一档，不再转给 \CTeX。取值：
%
% \begin{longtblr}[caption={\option{font-zh} 的取值}, label={tab:font-zh}]
%   {colspec={l X[2,l]}, width=\linewidth}
% \toprule
% 取值 & 字体 \\ \midrule
% \texttt{auto} & 按 \option{fontset} 决定，都没写则按平台猜 \\
% \texttt{windows} & 中易宋体、黑体、楷体、仿宋 \\
% \texttt{mac} & 华文宋体、黑体、楷体、仿宋 \\
% \texttt{ubuntu} & 思源（Noto CJK 版），名字跟 \CTeX\ 的同名 fontset 对齐 \\
% \texttt{fandol} & Fandol，\hologo{TeX} 发行版自带 \\
% \texttt{adobe} & Adobe Song/Hei/Kai/Fangsong Std \\
% \texttt{founder} & 方正书宋、黑体、楷体、仿宋 \\
% \texttt{hanyi} & 汉仪书宋二、中黑、楷体、仿宋 \\
% \texttt{none} & 不设置，用户自己 \cs{setCJKfamilyfont} 定义
%   \texttt{zhsong}、\texttt{zhhei}、\texttt{zhfs}、\texttt{zhkai} 四族 \\
% \bottomrule
% \end{longtblr}
%
% 模板早年提供过一个 \texttt{siyuan} fontset（\file{ctex-fontset-siyuan.def}），
% v3.2a 起去掉了。思源与 Noto CJK 是同一套字，用 \option{font-zh=ubuntu} 或
% \pkg{ctex} 的 \option{fontset=ubuntu} 即可。旧文档里写的
% \option{fontset=siyuan} 会被 \pkg{ctex} 报成找不到字体集，然后回退到平台默认
% 字体继续编译。
%
% \section{相邻文献序号的压缩}
% \label{sec:cite-compress}
% \option{bib} 族的 \option{compress-min} 管相邻连续的文献序号几个起压成
% 区间：
%
% \begin{longtblr}[caption={\option{compress-min} 的效果}, label={tab:compress}]
%   {colspec={l l l}}
% \toprule
% 取值 & \cs{cite}\{a,b\}（序号 1、2） & \cs{cite}\{a,b,c\}（序号 1、2、3） \\ \midrule
% \texttt{2} & \texttt{[1-2]}，默认 & \texttt{[1-3]} \\
% \texttt{3} & \texttt{[1,2]} & \texttt{[1-3]} \\
% \bottomrule
% \end{longtblr}
%
% 三个及以上连续的总归压成区间。不连续的序号本来就是逗号分隔，不受这个选项
% 影响。两条文献路线都认（见第~\ref{sec:bib-backend}~节），落点都是官方机制：
% 传统 \hologo{BibTeX} 一线转发给 \pkg{gbt7714.sty} 的 \texttt{mincompress}，
% biblatex 一线落到 \pkg{gb7714-2015} 的 \texttt{gbrefcompress} 计数器。
% 两边都不锁时机，什么时候设都行。
%
% \begin{latex}
% \hitsetup{bib={compress-min=3}}
% \end{latex}
%
% \section{参考文献的两条路线}
% \label{sec:bib-backend}
% 参考文献从 v3.2a 起分两条路线，\option{bib} 族的 \option{backend} 键选：
%
% \begin{longtblr}[caption={参考文献的两条路线}, label={tab:backend}]
%   {colspec={l X[1,l] l}, width=\linewidth}
% \toprule
% 取值 & 底下是什么 & 编译时跑的程序 \\ \midrule
% \texttt{bibtex} & \pkg{gbt7714}/\pkg{natbib} 加 \file{.bst}，默认 &
%   \texttt{bibtex} \\
% \texttt{biber} & \pkg{biblatex}，样式走 \pkg{biblatex-gb7714-2015} &
%   \texttt{biber} \\
% \bottomrule
% \end{longtblr}
%
% \texttt{biblatex} 当 \texttt{biber} 的别名也认。用 \texttt{latexmk} 编译不用
% 改别的——它看见 \file{.bcf} 就知道该跑 \texttt{biber}。不写这个键就是传统
% 一线；在导言区自己 \cs{usepackage}\marg{biblatex} 的，模板会侦测到，效果
% 等同 \texttt{biber}。
%
% \textbf{位置有讲究}：\pkg{biblatex} 必须赶在 \cs{begin}\marg{document} 之前
% 装载，所以 \option{backend}\texttt{=biber} 和 \option{struct} 里的
% \option{bibliography} 都要写在主文件的\emph{导言区}——demo 的
% \cs{hitsetup} 全在 \cs{begin}\marg{document} 之后，走 biber 一线时把这两组
% 往上挪：
%
% \begin{latex}
% \documentclass[…]{hithesis}
% \hitsetup[stage=final,lang=zh]{ struct = {
%   bibliography = {final/backmatter/reference} } }
% \hitsetup[stage=final]{ bib = { backend = biber } }
% \begin{document}
% \end{latex}
%
% 写晚了模板会一条条提醒，不会闷头编出错的结果。
%
% \option{bib} 族的通用键两条路线都认，名字对齐李泽平 gbt7714 v3 的官方
% 选项（kebab 拼写）：\option{uppercase-family}（\texttt{true} 著者姓名
% 全大写即国标排法，默认；\texttt{false} 保持 \file{.bib} 原样）、
% \option{max-bib-names}（文献表里最多列几位作者，\texttt{0} 不限）、
% \option{max-cite-names}（正文引用侧的同款，数字引用不出姓名，只在
% biblatex 一线对 \cs{textcite} 这类起作用）、\option{numbering-align}
% （文献表标号 \texttt{left}/\texttt{right}）、\option{compress-min}
% （见第~\ref{sec:cite-compress}~节）、\option{split-entry}（一条文献
% 可否跨页）、\option{punct}（著录标点三档，见下文换底座一段）、
% \option{comma-in-cite}（正文引用序号间的分隔：\texttt{half} 排
% \texttt{[3,5]} 默认、\texttt{half-with-space} 排 \texttt{[3,~5]}、
% \texttt{full} 排全角逗号）、\option{brackets-in-cite}（正文引用的方括号
% \texttt{half} 默认／\texttt{full} 排 ［3,5］）、\option{correct-punct}（\file{.bib} 数据里的
% 标点按 \option{punct} 档位清洗，默认开；biblatex 一线经 biber 的
% sourcemap 做同一套替换，姓名字段不动）、
% \option{sentence-case}（西文题名转不转句首大写。三个字段各一档：
% \texttt{sentence-case=\{title=true,booktitle=true,journal=false\}}；也可以
% \texttt{sentence-case=true|false|auto} 一次设三个。\texttt{auto} 是规范默认，
% 就是前面那一行的取值——期刊名不转，缩写形式本来就带自己的大小写。
% 名字取自两条后端线共有的说法，\file{.bst} 那边的控制字段叫
% \texttt{CTL\_sentence\_case\_title}，\pkg{biblatex} 那边的命令叫
% \cs{MakeSentenceCase}；不叫 \texttt{title-case}，那在西文排版里指每个实词
% 首字母大写，正好相反）、
% \option{show-mark}／\option{show-medium}（类型标 ［J］ 与载体标
% \texttt{/OL} 的显示开关，默认开）、\option{link-title}（题名加超链接，
% 默认关）、\option{show-url}／\option{show-doi}／\option{show-eprint}
% （默认开、开、关）、\option{show-note}（\file{.bib} 的 \texttt{note} 字段
% 排不排在条目末尾，默认关；成果页那种“SCI~收录，IDS~号为……”写进
% \texttt{note} 再打开它）、\option{show-sl-sn}（缺出版地或出版者时补不补
% ［出版地不详］/［出版者不详］，国标要补，默认 \texttt{true}，两条线
% 都可关）。除 \option{compress-min} 与 \option{split-entry} 外，其余在
% biblatex 一线是装载时换算成它的选项的，要么跟 \option{backend} 写进同一组
% \option{bib}\texttt{=\{...\}}，要么排在它前面，设晚了有提示。
%
% 后端专属的配置各归各的子组。传统一线换 \file{.bst} 写
% \option{bib}\texttt{=\{bibtex=\{style=gbt7714-2025-numeric\}\}}，值交给
% \cs{bibliographystyle}，\hologo{TeX} 树里任何 \file{.bst} 名都行，什么时候
% 设都行。biblatex 一线不开挑样式的口子——规范就是国标，样式锁死
% \texttt{gb7714-2015}；细调走 \option{biber}\texttt{=\{...\}}，整张键值表
% 原封并进 \pkg{biblatex} 的装载选项，\pkg{biblatex-gb7714-2015} 的几十个
% 选项（\texttt{gbpub}、\texttt{gbtype}、\texttt{gbannote}……）都从这儿过：
%
% \begin{latex}
% \hitsetup{ bib = { backend = biber ,
%                    biber = { gbpub = false , gbannote = true } } }
% \end{latex}
%
% 上标引用、\cs{inlinecite}、\cs{citeup}、全角序号、著录列表的版面都对齐到
% 传统一线。已知差别三处：著录条目的细排一边跟 \file{.bst}、一边跟
% \pkg{biblatex-gb7714-2015}，不逐字一致；传统一线默认把西文题名句首化
% （官方 \texttt{sentencecase=true}），biblatex 一线保持原样；
% \cs{inlinecite} 的标签前后没有传统一线那个中西文自动间距（gb 系样式特意
% 压掉的）；\option{punct} 在 biblatex 一线只落到方括号上（类型标、
% 访问日期），逗号冒号保持半角——2015 代的 biblatex 样式没有那半个旋钮
% （\texttt{gbpunctwidth} 是 2025 代才有的）。正文里残留的
% \cs{bibliographystyle} 会降级成一条提示，
% \cs{bibliography} 自动当 \cs{printbibliography} 排，老文档不用大改。
%
% 传统一线从 v3.2a 起只有一个 \file{hithesis.bst}，正文照抄李泽平官方
% \texttt{gbt7714-2015-numeric.bst}（v3.0.1）。原先深圳单独一个
% \file{hitszthesis.bst}，差异就是全角标点——标点如今是旋钮，不必再分文件：
% \option{bib} 族的 \option{punct} 三档，默认 \texttt{full}（不分语种
% 一律全角逗号冒号圆括号方括号、句点恒半角带空格，学校规范的示例与旧深圳版
% 都是这个排法）、\texttt{half}（一律半角，旧哈尔滨版的排法）、
% \texttt{by-lang}（按条目语言，中文全角西文半角）。全校区全学位同一个默认，
% 两条文献路线都认这个键，随设随生效。
%
% 对官方原版的补丁共五处（\file{src/hit-bst.dtx} 里 \texttt{HIT patch}
% 横幅标出）：\texttt{language} 认短码 \texttt{zh}/\texttt{en}（存量
% \file{.bib} 不用改）；默认档 \texttt{bylanguage}；句点恒半角带空格
% （官方全角档是 \texttt{．}，学校示例是 \texttt{.\ }）；方括号跟随全角档
% （类型标与访问日期排 ［Ｍ］ 这样的全角）；存量字段名 \texttt{citedate}
% 当 \texttt{urldate}、\texttt{typeoflit} 当 \texttt{entrytypeid} 的旧名
% （访问日期 ［1998-09-26］ 照旧显示）。其余默认照官方：西文姓氏全大写
% （\option{uppercase-family} 可关）、题名句首化、会议文献排出版地。
% 旧版丢出版地、\texttt{inbook} 排出悬空“\texttt{//}”的两个毛病随
% 换底座消失。写 \cs{bibliographystyle}\marg{hitszthesis} 的老文档请改成
% \texttt{hithesis}。
%
% \section{不再预载的 \pkg{tikz} 库}
% 模板过去会替用户 \cs{usetikzlibrary}\marg{calc}。v3.2a 起不再预载：类文件
% 自己已经不用它了（封面那个边框改用 \pkg{tikz} 核心的
% \texttt{[shift=\{…\}]} 写法），而预载它会逼着类文件在这一句上退出 expl3
% 语法（\cs{usetikzlibrary} 走裸 \cs{input}，不恢复 catcode）。
%
% 自己的插图里用到 \texttt{(\$…\$)} 坐标运算的，在导言区补一句
% \cs{usetikzlibrary}\marg{calc} 即可。
%
% \section{数学字体}
% \label{sec:math-fonts}
% 从 v3.2a 起数学公式走 \pkg{unicode-math}，数学字体由类选项 \option{math-font}
% 决定：
%
% \begin{longtblr}[caption={\option{math-font} 的取值}, label={tab:math-font}]
%   {colspec={l X[1,l]}, width=\linewidth}
% \toprule
% 取值 & 字体 \\ \midrule
% \texttt{auto} & 默认，等于 \texttt{xits} \\
% \texttt{xits} & XITS Math，Times 系，与正文字体最搭 \\
% \texttt{stix} & STIX Two Math \\
% \texttt{termes} & TeX Gyre Termes Math \\
% \texttt{newcm} & New Computer Modern Math \\
% \texttt{lm} & Latin Modern Math \\
% \texttt{libertinus} & Libertinus Math \\
% \texttt{newtx} & 老路子，8~位的 \pkg{newtxmath}，不走 \pkg{unicode-math} \\
% \texttt{none} & 不设置，用 \hologo{LaTeX} 默认的 CM \\
% \bottomrule
% \end{longtblr}
%
% 原来的 \option{newtxmath} 开关还在，默认值从 \texttt{true} 改成了
% \texttt{false}，写 \option{newtxmath=true} 等于 \option{math-font=newtx}。
%
% \option{math-style} 管数学排版风格：
%
% \begin{longtblr}[caption={\option{math-style} 的取值}, label={tab:math-style}]
%   {colspec={l X[1,l]}, width=\linewidth}
% \toprule
% 取值 & 含义 \\ \midrule
% \texttt{GB} & 默认。ISO 的基础上把 \cs{partial} 改直立 \\
% \texttt{ISO} & 按 ISO 80000-2，希腊字母一律倾斜 \\
% \texttt{TeX} & \hologo{TeX} 传统写法，大写希腊字母直立 \\
% \bottomrule
% \end{longtblr}
%
% 换引擎会影响已有文档的几处写法：
%
% \begin{itemize}
%   \item \cs{bm} 与 \cs{boldsymbol} 都接到了 \pkg{unicode-math} 的 \cs{symbf}
%     上。\pkg{unicode-math} 的符号是扩展 mathchar，\pkg{bm} 处理不了。
%   \item \cs{mathbf} 保持粗正体（接的是 \cs{symbfup}），与老版本一致。
%   \item \cs{mathrm}、\cs{mathit} 等仍从数学字体取字，而不是
%     \pkg{unicode-math} 默认的正文字体，这样 \cs{mathbf} 套希腊字母不会掉字。
%   \item \pkg{amssymb} 只在 \option{math-font=newtx} 与 \texttt{none} 下加载，
%     其余情况它提供的符号由数学字体自带。
% \end{itemize}
%
% \section{公式的物理量注释}
% \label{sec:eqdenote}
% \DescribeEnv{eqdenote}
% 规范要求公式中第一次出现的物理量代号给出注释，注释的转行与破折号后第一个字
% 对齐。\env{eqdenote} 把这套排版包好了，每条注释写一个 \cs{item}：
%
% \begin{latex}
% \begin{equation}\label{eq:1}
%   a = b + c
% \end{equation}
% \begin{eqdenote}
%   \item $a$ & 追踪飞行器与目标飞行器之间的相对位置矢量；
%   \item $b$ & 推力所产生的加速度。
% \end{eqdenote}
% \end{latex}
%
% 符号与说明之间用 \texttt{\&} 隔开，符号按原样写，要排数学就自己加
% \texttt{\$}。末条不用写 \cs{\textbackslash\textbackslash}。
%
% 可选参数换引导词，传空则不排引导词但保留等宽的空位：
%
% \begin{latex}
% \begin{eqdenote}[其中] ... \end{eqdenote}
% \begin{eqdenote}[]     ... \end{eqdenote}
% \end{latex}
%
% 表格底下是 \pkg{xltabular}（\pkg{longtable} 加 \texttt{X} 列），放不下会自动
% 转页接排，不用手动拆。要手动拆时把后半段另起一个 \env{eqdenote}\texttt{[]}，
% 两段的左缘仍然对齐。
%
% 引导词与破折号可以改：
%
% \begin{latex}
% \hitsetup{term={zh={eqdenote-label={式中}}, en={eqdenote-label={where}},
%                  mainmatter-eqdenote-dash={——}}}
% \end{latex}
%
% 中文文档用 \texttt{eqdenote-label-zh}，\option{lang=en} 时用
% \texttt{eqdenote-label-en}。
%
% \section{前文}
% \label{sec:titlepage}
% 前文内容是正文之前，含封面、摘要、目录、符号表、缩略语表。
% 封面信息提供两种配置方法：一是通过统一设置命 令 \cs{hitsetup}
% 通过\emph{key=value} 形式完成；二是每个信息利用命令独立设置， 其中命令的名字跟
% \emph{key} 相同。两种方式可以交叉使用，并按顺序执行（即后来的设置会覆
% 盖前面的）。以 \texttt{c} 开头的命令跟中文相关，\texttt{e}
% 开头则为对应的英文。
%
% \DescribeMacro{\hitsetup}
% \cs{hitsetup} 用法与常见 \emph{key=value} 命令相同，如下：
% \begin{latex}
% \hitsetup{
%   key1 = value1,
%   key2 = {a value, with comma},
% }
% % 可以多次调用
% \hitsetup{
%   key3 = value3,
%   key1 = value11, % 覆盖 value1
% }
% \end{latex}
%
% v3.2a 会在解析前去掉键值表里的空行，配置较长时可以用空行分组。
%
% 大多数命令的使用方法都是： \cs{command}\marg{arg}，例外者将具体指出。这些命令都
% 在示例文档的 \file{final/frontmatter/cover.tex} 中。
%
% \subsection{中英文两版的写法}
% \label{sec:setup-language}
% 中英文成对的字段写在 \option{zh} 与 \option{en} 两个分组里，字段名不带后缀：
%
% \begin{latex}
% \hitsetup{ info = {
%   zh = { title = {局部多孔质气体静压轴承关键技术的研究}, author = {于冬梅} },
%   en = { title = {Research on Key Technology of Aerostatic Bearing}, author = {Yu Dongmei} },
%   student-id = {1180800000},
% } }
% \end{latex}
%
% 示例文档使用这种写法。中英文字段各自集中在一个分组里，字段名不必重复写
% 语言后缀。不分语言的字段（\option{student-id}、
% \option{secret-level} 等）写在分组外面。
%
% \option{info} 与 \option{term} 两族都认这两个分组。
%
% 分组只是把字段名接上后缀再转发，所以带后缀的写法一样有效：
%
% \begin{latex}
% \hitsetup{ info = { title-zh = {...}, title-en = {...} } }
% \end{latex}
%
% 两者完全等价，混着写也行。带后缀的名字是键的本名，报错与改名提示里出现的是它，
% 下文各节也按本名列出，在分组里写的时候把后缀去掉就是了。
%
% 还有第三种写法，用斜杠把族名和字段名连起来，只设一两个字段时省一层括号：
%
% \begin{latex}
% \hitsetup{ info/zh/title = {...}, info/title-en = {...}, const/zh/school = {...} }
% \end{latex}
%
% 斜杠两边有没有空格都行。三种写法落到同一个键上，随便挑。
%
% 写错的分组会有提示：把不分语言的字段写进 \option{zh} 里，或者把一个根本不存在的
% 字段名写进去，都会说清是哪个键。斜杠写法同样。
%
% \subsection{密级}
% \label{sec:setup-secret}
% \DescribeMacro{secret-level}
% \DescribeMacro{classified-index-national}
% \DescribeMacro{classified-index-international}
% 定义秘密级别和国内国际索引号。
% \begin{latex}
% \hitsetup{
% secret-level={公开},
% classified-index-national={TM301.2},
% classified-index-international={62-5},
% }
% \end{latex}
%
% \subsection{论文标题}
% \myentry{论文标题}
% \DescribeMacro{title-zh}
% \DescribeMacro{title-en}
% \DescribeMacro{cover-title-zh}
% \DescribeMacro{title-line-one-zh}
% \DescribeMacro{title-line-two-zh}
% \DescribeMacro{subtitle-zh}
% \DescribeMacro{subtitle-en}
% 中英文标题，一个键管所有出现的地方。取值带两套断点：花括号分组给「题目栏是两条
% 固定宽度的线」那种版式分行（本科第二封面、博后第一封面、深圳博士中期封面），一组
% 一条线；\cs{\textbackslash\textbackslash} 给自由排的封面分行；方括号是副标题，
% 写了就排，不必再单独打开 \option{subtitle}。原创性声明的《\ldots》与 PDF 书签
% 取的是去掉断行的一整行。
%
% 副标题只出现在自由排的那几个封面上：硕博的第一、第二封面，哈尔滨本科第二封面，
% 以及英文封面。题目栏是两条固定线的三种封面——博后第一封面、威海本科第二封面、
% 深圳博士中期封面——版式里没有副标题这一栏，写了也不会排；报告整个没有副标题。
% 这是学校版式定的，不是模板漏了。
% \begin{latex}
% \hitsetup{
%   zh={
%     cover-title={封面中文题目可断行},
%     title={论文中文题目},
%     subtitle={论文中文副题目（如果有）},
%     title-line-one={本科生论文中文题目上部分},
%     title-line-two={本科生论文中文题目下部分},
%   },
%   en={
%     title={Thesis English Title},
%     subtitle={Thesis English Sub-Title (if necessary)},
%   },
% }
% \end{latex}
%
% \subsection{作者姓名}
% \myentry{作者姓名}
% \DescribeMacro{author-zh}
% \DescribeMacro{author-en}
% 作者姓名。
% \begin{latex}
% \hitsetup{
%   zh={author={中文姓名}},
%   en={author={Name in Pinyin}}
% }
% \end{latex}
%
% \subsection{申请学位名称}
% \label{sec:degree-level}
% \myentry{学科名称}
% \DescribeMacro{discipline-zh}
% \DescribeMacro{discipline-en}
% 按照入学的培养计划中学科自行填写，具体学科名称不是本文档范畴。
% 学位名称（博士/硕士/学士）由文档类的 \option{degree-level} 选项自动确定，无须手动设置。
%
% \begin{latex}
% \hitsetup{
%   zh={discipline={工学}},
%   en={discipline={Engineering}},
% }
% \end{latex}
%
% \subsection{院系名称}
% \myentry{院系名称}
% \DescribeMacro{affiliation-zh}
% \DescribeMacro{affiliation-en}
% 院系名称，同上，按照入学的培养计划中学科自行填写，具体院系名称不是本文档范畴。
% \begin{latex}
% \hitsetup{
%  zh={affiliation={机电工程学院}},
%  en={affiliation={\emultiline[t]{School of Mechatronics Engineering \\ Mechatronics Engineering}}},
% }
% \end{latex}
% \note[注意：]{个别学院英文名过长，使用以上方法自行换行。}
%
% \subsection{专业名称}
% \myentry{专业名称}
% \DescribeMacro{major-zh}
% \DescribeMacro{major-en}
% 专业名称，同上，按照入学的培养计划中学科自行填写，具体名称不是本文档范畴。
%
% 与上一节的 \option{discipline} 是两个字段：\option{discipline} 是学科门类，
% 封面上“工学博士”的“工学”取的是它；\option{major} 是专业，封面上“专业”
% 那一栏填的是它。两者各自独立，互不影响。
% \begin{latex}
% \hitsetup{
%  zh={major={机械制造及其自动化}},
%  en={major={Computer Science and Technology}},
% }
% \end{latex}
%
% \subsection{导师}
% \myentry{导师}
% \DescribeMacro{supervisor-zh}
% \DescribeMacro{supervisor-en}
% 直接导师。
% \begin{latex}
% \hitsetup{
%   zh={supervisor={导师~教授}},
%   en={supervisor={Supervisor}}
% }
% \end{latex}
%
% \myentry{副导师}
% \DescribeMacro{associate-supervisor-zh}
% \DescribeMacro{associate-supervisor-en}
% 副指导教师。
% \begin{latex}
% \hitsetup{
%   zh={associate-supervisor={副导师~副教授}},
%   en={associate-supervisor={2nd Boss}}
% }
% \end{latex}
%
% \myentry{联合导师}
% \DescribeMacro{co-supervisor-zh}
% \DescribeMacro{co-supervisor-en}
% 硕士生联合指导教师，博士生联合导师。
% \begin{latex}
% \hitsetup{
%   zh={co-supervisor={联合导师~教授}},
%   en={co-supervisor={3rd Boss}}
% }
% \end{latex}
%
% \subsection{成文日期}
% \myentry{日期}
% \DescribeOption{date-zh}
% \DescribeOption{date-en}
% \textbf{日期填数字，排法交给模板。} 不写就取编译当天（东八区）的年月；
% 要指定就写 \texttt{yyyy}、\texttt{yyyy-mm} 或 \texttt{yyyy-mm-dd}：
% \begin{latex}
% \hitsetup{ info = {
%   zh = { date = {2026-06} },
%   en = { date = {2026-06} },
% } }
% \end{latex}
% 同一个数字，中文封面排成“2026 年 6 月”，英文封面排成“June, 2026”。
% 月和日补不补零都认，\texttt{2026-7-7} 与 \texttt{2026-07-07} 一样。
%
% \textbf{写到哪级排到哪级。} 写年月排年月，写年月日排年月日，只写年就只排年，
% 模板不截也不催。规范四处落点要的都是年月，所以默认值只给到月；深圳本科的
% 开题中期封面要年月日，那一档的默认值单独给全。
%
% \textbf{写成别的样子就自己负责。} 值不合 ISO 就原样排出来，另在日志里说一句。
% 想排“二〇二六年六月”这种汉字数字，写
% \cs{zhdigits}\marg{2026}\texttt{年}\cs{zhnumber}\marg{6}\texttt{月}——
% 年份要用 \cs{zhdigits} 逐位念，\cs{zhnumber}\marg{2026} 会念成“二千零二十六”。
%
% \myentry{博士后日期}
% \DescribeMacro{research-start-date-zh}
% \DescribeMacro{research-end-date-zh}
% \DescribeMacro{report-submit-date-zh}
% \DescribeMacro{work-finish-date-zh}
% 博士后出站报告封面需要填写进站、出站、提交、完成等多个日期，对应命令分别为
% \cs{research-start-date-zh}、\cs{research-end-date-zh}、\cs{report-submit-date-zh}、\cs{work-finish-date-zh}。
% 其余学生类型的封面无须填写。
%
% \subsection{学生类型}
% \myentry{学生类型}
% \DescribeMacro{student-type-zh}
% \DescribeMacro{student-type-en}
% 非全日制教育申请学位者
%（同等学力人员）、（工程硕士）、（工商管理硕士）、
%（高级管理人员工商管理硕士）、（公共管理硕士）、（中职教师）、（高校教师）等，
% 具体要求按照入学的培养计划中学科自行填写，具体名称不是本文档范畴。
% \begin{latex}
% \hitsetup{
%   zh={degree-type={同等学力人员}},
%   en={degree-type={Master of Art}},
% }
% \end{latex}
%
% \subsection{学号}
% \myentry{学号}
% \DescribeMacro{student-id}
% \DescribeMacro{student-number}
% 学号，具体要求按照入学的培养计划中学科自行填写。
% 两个名字都对，写哪个都不提示。
% \begin{latex}
% \hitsetup{
%   student-id={9527},
% }
% \end{latex}
%
% \subsection{摘要}
% \myentry{摘要正文}
% \DescribeEnv{abstract-zh}
% \DescribeEnv{abstract-en}
% \note[说明：]{摘要正文只能用环境命令的形式，不支持 \cs{hitsetup}。}
%
% \begin{latex}
% \begin{abstract-zh}
%  摘要请写在这里...
% \end{abstract-zh}
%
% \begin{abstract-en}
%  Here comes the abstract in English...
% \end{abstract-en}
% \end{latex}
%
% \myentry{关键词}
% \DescribeMacro{keywords-zh}
% \DescribeMacro{keywords-en}
% 关键词用英文逗号分割写入相应的命令中，模板会解析各关键词并生成符合不同论文格式
% 要求的关键词格式。
% \begin{latex}
% \hitsetup{
%   zh={keywords={关键词 1, 关键词 2}},
%   en={keywords={keyword 1, keyword 2}}
% }
% \end{latex}
%
% \subsection{符号表}
% \DescribeEnv{symbols}
% 符号表环境，单独放一个文件，在 \option{struct} 里写
% \option{list-of-symbols}\texttt{=\{final/frontmatter/symbols\}}。
% 一行一个符号，\cs{item}\oarg{符号}\meta{说明}：
% \begin{latex}
% \begin{symbols}
% \item[$\rho$]        密度，$\mathrm{kg}/\mathrm{m}^{3}$
% \item[$\mathbb{R}$]  实数集
% \item[Hz]            频率单位赫[兹]，$\mathrm{s}^{-1}$
% \end{symbols}
% \end{latex}
% 顺序完全由自己定，模板不排序。\PGR\ 没有规定符号表怎么排；投国际期刊常见的
% “拉丁字母、希腊字母、数字、特殊符号”那种分组顺序，手写也排得出来。
%
% 可选参数收键值：\option{labelwidth} 改第一列宽度（默认 2.5\,cm），
% \option{subheading} 给这一段加个小标题（只在一页里排两段时用得着）。
% 不含等号的裸值当宽度，所以 v3.1e 的写法照样认：
% \begin{latex}
% \begin{symbols}[3cm]
% \begin{symbols}[labelwidth=3cm, subheading=符号]
% \end{latex}
%
% \DescribeEnv{symbols*}
% \PGR\ 允许把符号分成几张表、各有各的列，列表排不出那种版式。要的话环境名加个
% 星号，标题和表号照旧由环境给，正文完全由自己写。
% \begin{latex}
% \begin{symbols*}
% \begin{longtable}{cccc}%
% \toprule 量的名称 & 单位名称 & 单位符号 & 其它表示实例\\ \midrule
% \endfirsthead
% \multicolumn{4}{r}{表~\thetable（续表）}\vspace{0.5em}\\
% \toprule 量的名称 & 单位名称 & 单位符号 & 其它表示实例\\ \midrule
% \endhead
% \bottomrule
% \endfoot
% 频率 & 赫[兹] & Hz & s$^{-1}$\\
% 力 & 牛[顿] & N & kg$\cdot$m/s$^2$\\
% \end{longtable}
% \end{symbols*}
% \end{latex}
%
% \DescribeEnv{denotation}
% v3.1e 的环境名，只供老文档：排标题、重置表号，正文完全由自己写，行为与 v3.1e
% 一样。新写法一律用 \env{symbols}。
%
% \subsection{缩略语表}
% \label{sec:abbr}
% \DescribeMacro{\abbr}
% \DescribeMacro{\abbrs}
% 缩略语的词条填在 \option{abbr} 族里，键名就是缩写：
% \begin{latex}
% \hitsetup{ abbr = {
%   SCNA  = {zh=体细胞拷贝数变异, en=Somatic copy number alteration},
%   TSSBP = {zh=树结构折筷过程, en=Tree-structured Stick-breaking process},
%   SVM   = {Support Vector Machine},
% } }
% \end{latex}
% 值收 \option{zh} 与 \option{en}；只有一种语言时写裸值，含汉字归中文、不含归英文。
% \option{label} 改标签，不写就取缩写本身。
%
% 正文里写 \cs{abbr}\marg{缩写}，首次出现全写、之后只出缩写；\cs{abbrs} 是复数；
% 带星号是强制全写且不改“用过”的状态，摘要或新一章想再展开一次时用。
% \PGR\ 要求“文中第一次出现的缩写词应该用括号注明英文原词”，三种取值各一种排法：
% \begin{longtblr}[caption={缩写首次出现的三种排法}, label={tab:abbr}]
%   {colspec={l X[1,l] X[1,l]}, width=\linewidth}
% \toprule
% 词条 & 中文版首次 & 英文版首次 \\ \midrule
% \option{zh}+\option{en} & 体细胞拷贝数变异（Somatic copy number alteration，SCNA） & Somatic copy number alteration (SCNA) \\
% 只有 \option{zh}        & 树结构折筷过程（TSSBP） & 树结构折筷过程 (TSSBP) \\
% 只有 \option{en}        & SVM（Support Vector Machine） & Support Vector Machine (SVM) \\
% \bottomrule
% \end{longtblr}
% 摘要排完会把“用过”的状态清一次，正文第一次再出现仍全写一遍——读者未必先读
% 摘要，\PGR\ 说的“文中第一次出现”指的也是正文。不想清就写
% \option{abbr}/\option{reset-after-abstract}\texttt{=false}。
%
% 括号与逗号取全角还是半角由 \option{abbr}/\option{punct} 定：\texttt{auto}（默认）
% 按语种走，中文版全角、英文版半角；中文数学类论文行文用半角标点的写 \texttt{half}。
%
% \DescribeMacro{\listofabbreviations}
% 那一页由类自己生成，在 \option{struct} 里写
% \option{list-of-abbreviations}\texttt{=true}。版式与符号表同一套两列悬挂：
% 缩写一列顶格，说明一列从固定位置排到右边距、长了折行悬挂；中文版说明是
% “中文全称（英文全称）”合成一格，只有一种语言就排那一种，英文版只排英文
% 全称（没有就中文）。都按缩写的字母序排。\option{labelwidth} 键改第一列
% 宽度，与 \env{symbols} 同名同义，合并成一页时两段同宽说明列就对齐。
%
% 手排那一页写 \cs{listofabbreviations}，与 \cs{listoffigures} 一族同名同姓。
% 带星号是不排章标题，接在符号表那一页下面用。可选参数收 \option{subheading}，
% 给的是那一段的小标题。
%
% \DescribeMacro{\hitindex}
% 缩写用一次就顺带往中英两个索引各登记一条，索引形式直接取 \option{zh} 与
% \option{en}。不想登记就写 \option{abbr}/\option{indexing}\texttt{=false}，
% 单处不想登记写 \cs{abbr}\oarg{indexing=false}\marg{缩写}。
%
% 不缩写的词用 \cs{hitindex}：
% \begin{latex}
% 北\hitindex{zh=乔峰, en=Qiao Feng}南慕容
% \end{latex}
% 排出当前语种的那一份，两个索引都登记。裸值按有没有汉字归边。带星号只登记不排，
% 正文里写的是“乔帮主”而索引条目要“乔峰”时用。
%
% 中文索引要一个拼音排序键（\pkg{makeindex} 不认汉字），模板查表推首字的拼音并把
% 声调转成数字，“体细胞拷贝数变异”得到 \texttt{ti3}。碰上多音字或库里查不到会
% 提示一次，那时候手写 \option{pinyin}\texttt{=ti3} 覆盖。
%
% \subsection{符号表与缩略语表排在同一页}
% 不少论文只列一张“本文符号及缩写”。这时在 \option{struct} 里写第三个部件：
% \begin{latex}
% list-of-symbols-and-abbreviations = {final/frontmatter/symbols},
% \end{latex}
% 一个章标题“符号及缩略语”，页内两段各带小标题，符号在上、缩略语在下。写它就是
% 一页，写前两个部件就是两页。
%
% \subsection{目录}
% 目录不需要用户干预，自动生成，具体命令已经写在 \file{final.tex} 中。
%
% \section{正文}
%
% \subsection{标题层级}
% \label{sec:heading-level}
% 终稿的标题是章、节、条、款四级，写 \cs{chapter}、\cs{section}、\cs{subsection}、
% \cs{subsubsection}。报告没有章，顶级是节，往下是条、款。
%
% 报告的正文从 \cs{section} 写起或从 \cs{chapter} 写起都行：\emph{正文里第一个带编号的
% 标题命令是哪一级，那一级就排成节}，其余依次往下。所以按终稿结构写的一份正文，
% 改个 \option{stage} 就能排成报告；从 \cs{section} 写起的老习惯也照旧。
% \begin{longtblr}[caption={报告与终稿的标题层级}, label={tab:heading-level}]
%   {colspec={l X[1,l] X[1,l]}, width=\linewidth}
% \toprule
% 正文从哪一级写起 & 报告排成 & 终稿排成 \\ \midrule
% \cs{chapter} & \cs{chapter} 是节 1，\cs{section} 是条 1.1，\cs{subsection} 是款 1.1.1 & 章、节、条、款 \\
% \cs{section} & \cs{section} 是节 1，\cs{subsection} 是条 1.1，\cs{subsubsection} 是款 1.1.1 & 终稿必须有章，要先写 \cs{chapter} \\
% \bottomrule
% \end{longtblr}
% 终稿不做这种换算，从 \cs{section} 写起的终稿不会被当成章。
%
% 规范没有规定的层级不给兜底，下面几种写法直接报错：
% \begin{itemize}
% \item 报告从 \cs{section} 写起，后面又写了 \cs{chapter}。
% \item 报告正文的第一个编号标题是 \cs{subsection} 或 \cs{subsubsection}，前面没有顶级。
% \item 落在款以下：\cs{paragraph}、\cs{subparagraph}，以及报告从 \cs{chapter} 写起时的
%   \cs{subsubsection}。规范的标题只到款，再往下请写列表（\env{enumerate}、\env{itemize}）。
%   旧版面（v3.1e 兼容）不拦 \cs{paragraph}。
% \end{itemize}
%
% \subsection{图和表}
% \hit\ 博士毕业论文要求使用中英双语图题、表题。写法是
% \cs{caption}\marg{中文题}\oarg{英文题}，英文题作可选参数跟在中文题后面，
% 类会另起一行排。两个题之间的距离\PGR\ 没有规定，由类给一个固定值，不用自己调。
% 关于图题\PGR\ 和\UGR\ 只规定了居中，并没有规定居中对齐。然而评审老师很多喜欢居
% 中且居中对齐。模板默认选项是居中但不居中对齐，如果需要居中对齐，那么需要在
% \file{final.tex}的文档类选项中设置选项capcenterlast=true。详细方法见前文的介绍
% 。
% \begin{hitrgu}
% 每个图均应有图题（由图序和图名组成），图题不宜有标点符号，图名在图序之后空1个
% 半角字符排写。图序按章编排，如第1章第一个插图的图号为“图1-1”。图题置于图下，硕
% 士论文只用中文，博士论文用中、英两种文字，居中书写，中文在上，要求中文用宋体5
% 号字，英文用Times New Roman 5号字。有图注或其它说明时应置于图题之上。引用图应
% 注明出处，在图题右上角加引用文献号。图中若有分图时，分图题置于分图之下或图题之
% 下，可以只用中文书写，分图号用a)、b)等表示。图中各部分说明应采用中文（引用的外
% 文图除外）或数字符号，各项文字说明置于图题之上（有分图时，置于分图题之上）。图
% 中文字用宋体、Times New Roman字体，字号尽量采用5号字（当字数较多时可用小5号字
% ，以清晰表达为原则，但在一个插图内字号要统一）。同一图内使用文字应统一。图表中
% 物理量、符号用斜体。
% \end{hitrgu}
% \DescribeMacro{\caption}
% 双语题注写成 \cs{caption}\marg{中文题}\oarg{英文题}：
%
% \begin{latex}
% \begin{figure}
%   \includegraphics{golfer}
%   \caption{打高尔夫球的人}[The person playing golf]
% \end{figure}
% \end{latex}
%
% 不给尾随的 \oarg{英文题} 就只排中文题。目录用的短标题仍是 \cs{caption} 原本那个
% 前置可选参数，两者位置不同不会混：\cs{caption}\oarg{短标题}\marg{中文题}\oarg{英文题}。
%
% 英文图表名（\texttt{Fig.}、\texttt{Table}）不用每次手写，取自 \option{term} 族的
% \option{figure-prefix-en} 与 \option{table-prefix-en}。
%
% \DescribeMacro{\figurenote}
% 图注写成 \cs{figurenote}\marg{文字}，放在 \cs{caption} \emph{前面}——规范要求
% 图注置于图题之上：
% \begin{latex}
% \begin{figure}
%   \includegraphics{golfer}
%   \figurenote{图中箭头表示演化方向。}
%   \caption{打高尔夫球的人}[The person playing golf]
% \end{figure}
% \end{latex}
% 字号行距由类给（\option{styles} 族的 \option{caption-note}，与题注同一档），
% 不用自己写 \cs{wuhao} 与 \cs{vspace}。
%
% 题注本身也一样：\option{caption-figure}、\option{caption-table}、
% \option{caption-sub} 三个目标定图题、表题与分图题的字号行距，默认按写作指南
% 排宋体五号、行距 1.25。要改写 \cs{hitsetup}\marg{format~=~\{~caption-figure~=~\{~\ldots~\}~\}}，
% \option{caption} 这个 meta 目标一次改图题与表题两处。
%
% \env{table} 浮动体里的默认字号也是五号（规范：“表中文字用宋体~5~号字”），
% 长表格同样，不必自己在表格前写 \cs{wuhao}。
%
% \DescribeMacro{\hitsetup}
% \option{styles/bilingual-caption} 三态：\texttt{auto}（默认）按学位——博士与
% 博士后排双语，本科与硕士只排中文题，写了尾随的英文题也不排，规范原话是
% “硕士学位论文只用中文，博士学位论文用中、英两种文字”。硕士确要双语写
% \texttt{true} 强制打开；\texttt{false} 一概不排，同一份稿子投不同用途时不必
% 逐条删。开着而某条题注没给英文题，那一条按单语排，不报错也不留空。
%
% 中文题太长、插图索引里放不下时，用 \cs{caption} 原本那个前置可选参数放短标题：
% \begin{latex}
% \caption[短标题]{很长的中文题注……}[The person playing golf]
% \end{latex}
%
% 英文图表名取 \option{term} 族的 \option{figure-prefix-en} 与 \option{table-prefix-en}
% （默认 \texttt{Fig.} 与 \texttt{Table}，都已经带好负细空格，用来抵掉
% \cs{fnum@figure} 插的那个空格）。要换成别的名字：
% \begin{latex}
% \hitsetup{term={caption-figure-prefix-en={Figure~}}}
% \end{latex}
%
% \DescribeMacro{\bicaption}
% \cs{bicaption} 是 v3.2a 之前的写法，仍然可用但每种形式各提示一次，v4 移除。
% 三种形式都接得住：\pkg{ccaption} 原本的五参数
% \cs{bicaption}\oarg{标签}\marg{索引短标题}\marg{中文题}\marg{英文图表名}\marg{英文题}，
% 以及本模板之前的三参数与四参数。新写的文档请用 \cs{caption}。
%
% \DescribeMacro{\longbionenumcaption}
% 长表格的双语题注原先要写 \cs{longbionenumcaption}，七个参数依次是中文短标题、
% 中文题、英文表名、英文短标题、英文题、中英标题间距、题注与表格的间距。现在
% \env{longtable} 里也可以直接写 \cs{caption}\marg{中文题}\oarg{英文题}，跟浮动体
% 一样；只要中文题就不写那个可选参数。后两个间距参数由类里统一给，不必再自己算。
% \cs{longbionenumcaption} 仍然可用，用到提示一次，v4 移除。
%
% 题注与图表之间的空白由类统一给，正文里不用再写 \cs{vspace}\marg{0.5em}。表题在上
% 就放在题注下面，图题在下就放在题注上面。浮动体与正文、浮动体与浮动体、页顶页底、
% 整页浮动，这四种位置的间距现在取同一个值，所以同一份文档里几张图表的前后空白是
% 一样宽的。旧文档里那些手写的 \cs{vspace} 会叠加上去，建议删掉。
%
% 单双语图题的方法如下，注释中说明。
% \begin{latex}
% \begin{figure}[htpb]
% \centering
% \includegraphics[width = 0.4\textwidth]{golfer}
% \caption[golfer1]{注意图中文字尽量用五号字
% }[The person playing golf]
% \end{figure}
% \end{latex}
% 单张单图题的格式如下，
% \begin{latex}
% \begin{figure}[h]
% \centering
% \includegraphics[width = 0.4\textwidth]{golfer}
% \caption{注意图中文字字号尽量用五号字}
% \end{figure}
% \end{latex}
% 并排图例。
% \begin{latex}
% \begin{figure}[htbp]
% \centering
% \begin{minipage}{0.4\textwidth}
% \centering
% \includegraphics[width=\textwidth]{golfer}
% \caption[golfer2]{打高尔夫球的人}[The person playing golf]
% \end{minipage}
% \begin{minipage}{0.4\textwidth}
% \centering
% \includegraphics[width=\textwidth]{golfer}
% \caption[golfer3]{打高尔夫球的人}[The person playing golf]
% \end{minipage}
% \end{figure}
% \end{latex}
% 子图图例。子图题只要一种语言时直接写在可选参数里：
% \begin{latex}
% \begin{figure}[htbp]
% \centering
% \subfigure[打高尔夫球的人~1]{\label{golfer41}\includegraphics[width=0.4\textwidth]{golfer}}
% \subfigure[打高尔夫球的人~2]{\label{golfer42}\includegraphics[width=0.4\textwidth]{golfer}}
% \caption[golfer4]{打高尔夫球的人}[The person playing golf]
% \end{figure}
% \end{latex}
%
% \DescribeMacro{\bisubcaption}
% 要中英文两行子图题，把 \cs{bisubcaption}\marg{中文}\marg{英文} 写进可选参数：
% \begin{latex}
% \begin{figure}[htbp]
% \centering
% \subfigure[\bisubcaption{打高尔夫球的人~1}{The person playing golf}]{%
%   \label{golfer41}\includegraphics[width=0.4\textwidth]{golfer}}
% \subfigure[\bisubcaption{打高尔夫球的人~2}{The person playing golf}]{%
%   \label{golfer42}\includegraphics[width=0.4\textwidth]{golfer}}
% \caption[golfer4]{打高尔夫球的人}[The person playing golf]
% \end{figure}
% \end{latex}
%
% v3.2a 之前子图由 \pkg{subfigure} 实现，双语子图题只能靠嵌套两层 \cs{subfigure}
% 再拿 \cs{addtocounter} 把计数器凑回来，标签也得用一个空的 \cs{subfigure} 单独占：
% \begin{latex}
% \subfigure{\label{golfer41}}\addtocounter{subfigure}{-2}
% \subfigure[The person playing golf]{\subfigure[打高尔夫球的人~1]{\includegraphics{golfer}}}
% \end{latex}
% \pkg{subfigure} 二十多年没人维护，与 \pkg{cleveref}、\pkg{hyperref} 配合有问题，
% 现已换成 \pkg{subcaption}。上面这种旧写法仍然可用，模板会自动转成新结构并在日志里
% 提示，v4（预计~2027~年发布）移除。排出来的版面与旧版一致。
%
% \pkg{subfigure} 专有的开关与 \cs{listsubcaptions} 已经失效，用到会在日志里报出来
% 并指出替代写法。都能在 \pkg{caption} 的 \cs{captionsetup}\oarg{sub} 里找到对应：
%
% \begin{longtblr}[caption={\pkg{subfigure} 旧写法与 \pkg{subcaption} 新写法对照}, label={tab:subfig}]
%   {colspec={l X[1,l]}, width=\linewidth}
% \toprule
% 旧写法 & 新写法 \\ \midrule
% \cs{subcaphangtrue}        & \cs{captionsetup}[sub]\{format=hang\} \\
% \cs{subcapcentertrue}      & \cs{captionsetup}[sub]\{justification=centering\} \\
% \cs{subcapcenterlasttrue}  & \cs{captionsetup}[sub]\{justification=centerlast\} \\
% \cs{subcapnoonelinetrue}   & \cs{captionsetup}[sub]\{singlelinecheck=false\} \\
% \cs{subcapraggedrighttrue} & \cs{captionsetup}[sub]\{justification=raggedright\} \\
% \cs{listsubcaptions}       & \cs{captionsetup}[sub]\{list=true\}（不自动转） \\
% \bottomrule
% \end{longtblr}
%
% 把 \texttt{true} 换成 \texttt{false} 就是关掉，转成对应键的默认值。要注意
% \option{justification} 只有一个值，\cs{subcapcenter}、\cs{subcapcenterlast}、
% \cs{subcapraggedright} 三个开关都写它，后设的会盖掉先设的；\pkg{subfigure} 那边
% 是三个独立开关。
%
% \cs{listsubcaptions} 不自动转：它原先是在指定位置把攒下来的子图题一次性写进插图
% 索引，而 \option{list} 是每条子图题排出来时就写，等调用到 \cs{listsubcaptions}
% 已经晚了。用到只会提示，请自己在导言区写
% \cs{captionsetup}\oarg{sub}\marg{list=true}。
%
% \cs{subfigtopskip} 这一批长度仍然有效。
%
% 一张图分两页排时，后一页那半张写成续图：\cs{ContinuedFloat} 让子图接着上一张
% 编号，\cs{caption*} 排不编号的题注，只印写进去的正文。
% \begin{latex}
% \begin{figure}
%   \ContinuedFloat
%   \centering
%   \subfigure[气体动静压轴承]{\includegraphics[width=6.7cm]{hybrid}}%
%   \hspace{1cm}%
%   \subfigure[气体压膜轴承]{\includegraphics[width=7.2cm]{squeeze-film}}
%   \caption*{图~\thefigure（续图）}
% \end{figure}
% \end{latex}
% 子图之间的行尾记得用 \verb|%| 吃掉空格，两张图加上间距刚好卡着版心宽时，
% 多出来的词间空格会把第二张挤到下一行。
%
% \textbf{表格推荐用 \pkg{tabularray} 排。}类已经装好它连同 booktabs 库，
% 导言区什么都不用写。新老写法共用一套值，类都接好了：\cs{toprule} 一族画出来的线一样粗
% （1.5\,bp／1\,bp／0.5\,bp），单元格左右边距同为 5.4\,bp，浮动体里字号
% 同样压到五号，表题同一条链。用哪种写法版面一致；新文档建议用
% \pkg{tabularray}——列格式集中写在 \texttt{colspec} 一处，跨页长表不用
% 手写续表头。
%
% 新写法（\env{booktabs} 环境由 \cs{UseTblrLibrary}\marg{booktabs} 提供，
% \cs{toprule} 一族在其中可用）：
% \begin{latex}
% \begin{table}[htbp]
% \caption[table1]{符合研究生院绘图规范的表格}[Table in agreement of the standard from graduate school]
% \centering
% \begin{booktabs}{colspec={ccccc}}
% \toprule
% $D$(in) & $P_u$(lbs) & $u_u$(in) & $\beta$ & $G_f$(psi.in)\\
% \midrule
%  5 & 269.8 & 0.000674 & 1.79 & 0.04089\\
% 10 & 421.0 & 0.001035 & 3.59 & 0.04089\\
% 20 & 640.2 & 0.001565 & 7.18 & 0.04089\\
% \bottomrule
% \end{booktabs}
% \end{table}
% \end{latex}
% 同一张表的老写法（v3.1e 一直的写法，仍然可用）只是把表体换成
% \env{tabular}：
% \begin{latex}
% \begin{tabular}{ccccc}
% \toprule
% $D$(in) & $P_u$(lbs) & $u_u$(in) & $\beta$ & $G_f$(psi.in)\\
% \midrule
%  5 & 269.8 & 0.000674 & 1.79 & 0.04089\\
% \bottomrule
% \end{tabular}
% \end{latex}
% 两种写法都不必再写 \cs{wuhao}：\env{table} 浮动体里默认就是五号。
%
% \textbf{长表格。}长表格不是浮动体，不会自动调整位置，排在写的地方。
% 推荐用 \env{longtblr}：
% \begin{latex}
% \begin{longtblr}[caption={中国省级行政单位一览}, label={table2}]
%   {colspec={ccc}}
% \toprule 名称 & 简称 & 省会或首府 \\ \midrule
% 北京市 & 京 & 北京\\
% 天津市 & 津 & 天津\\
% 河北省 & 冀 & 石家庄市\\
% \bottomrule
% \end{longtblr}
% \end{latex}
% 续页自动重复表头（\texttt{rowhead} 默认 1，写 0 关掉）；续页右上角的
% “表~x-x（续表）”由类排（规范原话：“表右上角注明编号，编号后加
% ‘（续表）’，表题可省略”），不用手写 \cs{endfirsthead} 那一套。字号行距
% 与浮动表同规矩。博士论文的英文表题写在 \texttt{caption-en} 外参里：
% \begin{latex}
% \begin{longtblr}[caption={中文表题}, caption-en={English title},
%   label={tbl}]{colspec={ccc}}
% \end{latex}
% 英文行的表名、表号、间隔与浮动表 \cs{caption}\marg{中文题}\oarg{英文题}
% 的第二行同一条链，排不排也听 \option{styles}/\option{bilingual-caption}
% 同一个开关。
%
% 老写法 \pkg{longtable} 仍然可用，双语题注
% \cs{caption}\marg{中文题}\oarg{英文题}，续表头自己写：
%
% \textbf{字号与行距不用自己设。} 类里已经把长表格默认压到五号，行距取版心的
% 网格行距——Word 里表格数据行的行高就是文档网格的行距，范例那张表量出来一行
% 19.18\,bp，与 \texttt{docGrid} 写的 \texttt{linePitch=383}（19.15\,bp）对得上。
% v3.1e 教的 \cs{ltfontsize}\marg{\cs{wuhao}\oarg{1.667}} 手调不必再写；
% 确实要另一个行距时 \cs{ltfontsize} 仍然管用。
% \begin{latex}
% \begin{longtable}{ccc}
% \caption{中国省级行政单位一览}[Overview of the provincial administrative
% unit of China]\label{table2}\\
% %\caption{中国省级行政单位一览}\label{table2}\\ %不写可选参数就是单语标题
% \toprule 名称 & 简称 & 省会或首府  \\ \midrule
% \endfirsthead
% \multicolumn{3}{r}{表~\thetable（续表）}\vspace{0.5em}\\
% \toprule 名称 & 简称 & 省会或首府  \\ \midrule
% \endhead
% \bottomrule
% \endfoot
% 北京市 & 京 & 北京\\
% 天津市 & 津 & 天津\\
% 河北省 & 冀 & 石家庄市\\
% 山西省 & 晋 & 太原市\\
% 内蒙古自治区 & 蒙 & 呼和浩特市\\
% 辽宁省 & 辽 & 沈阳市\\
% 吉林省 & 吉 & 长春市\\
% 黑龙江省 & 黑 & 哈尔滨市\\
% 上海市 & 沪/申 & 上海\\
% 江苏省 & 苏 & 南京市\\
% 浙江省 & 浙 & 杭州市\\
% 安徽省 & 皖 & 合肥市\\
% 福建省 & 闽 & 福州市\\
% 江西省 & 赣 & 南昌市\\
% 山东省 & 鲁 & 济南市\\
% 河南省 & 豫 & 郑州市\\
% 湖北省 & 鄂 & 武汉市\\
% 湖南省 & 湘 & 长沙市\\
% 广东省 & 粤 & 广州市\\
% 广西壮族自治区 & 桂 & 南宁市\\
% 海南省 & 琼 & 海口市\\
% 重庆市 & 渝 & 重庆\\
% 四川省 & 川/蜀 & 成都市\\
% 贵州省 & 黔/贵 & 贵阳市\\
% 云南省 & 云/滇 & 昆明市\\
% 西藏自治区 & 藏 & 拉萨市\\
% 陕西省 & 陕/秦 & 西安市\\
% 甘肃省 & 甘/陇 & 兰州市\\
% 青海省 & 青 & 西宁市\\
% 宁夏回族自治区 & 宁 & 银川市\\
% 新疆维吾尔自治区 & 新 & 乌鲁木齐市\\
% 香港特别行政区 & 港 & 香港\\
% 澳门特别行政区 & 澳 & 澳门\\
% 台湾省 & 台 & 台北市\\
% \end{longtable}
% \end{latex}
% \subsection{公式}
% 公式本身与正常用法一致。
%
% \DescribeOption{math/numbering-parenthesis-style}
% 公式编号两边的括号取全角还是半角。取 \texttt{full} 是全角（默认），
% \texttt{half} 是半角。
% \begin{latex}
% \hitsetup{math={numbering-parenthesis-style=half}}
% \end{latex}
% \cs{eqref} 跟着一起变，正文里的引用与公式旁的编号不会一个全角一个半角。
%
% \option{lang} 取 \texttt{en} 时这个键不起作用，一律半角。英文正文里夹一对全角
% 括号，宽度和字距都不对。
%
% 键在导言区任何位置写都行，甚至正文里改也跟得上：括号是排的时候才判断的，不是
% 设键时定死的。
% \subsection{数学环境}
% \label{sec:math}
% \hithesis\ 定义了常用的数学环境：
%
% \begin{center}
% \begin{tabular}{*{7}{l}}\toprule
%   axiom & theorem & definition & proposition & lemma & conjecture &\\
%   公理 & 定理 & 定义 & 命题 & 引理 & 猜想 &\\\midrule
%   proof & corollary & example & exercise & assumption & remark & problem \\
%   证明 & 推论 & 例子& 练习 & 假设 & 注释 & 问题\\\bottomrule
% \end{tabular}
% \end{center}
%
% 比如：
% \begin{latex}
% \begin{definition}
%   道千乘之国，敬事而信，节用而爱人，使民以时。
% \end{definition}
% \end{latex}
% 产生（自动编号）：
% \medskip
%
% \noindent\framebox[\linewidth][l]{{\heiti 定义~1.1~~~} % {道千乘之国，敬事而信，节用而爱人，使民以时。}}
%
% \smallskip
% 列举出来的数学环境毕竟是有限的，如果想用\emph{胡说}这样的数学环境，那么可以定义：
% \begin{latex}
% \newtheorem{nonsense}{胡说}[chapter]
% \end{latex}
%
% 然后这样使用：
% \begin{latex}
% \begin{nonsense}
%   契丹武士要来中原夺武林秘笈。—— 慕容博
% \end{nonsense}
% \end{latex}
% 产生（自动编号）：
%
% \medskip
% \noindent\framebox[\linewidth][l]{{\heiti 胡说~1.1~~~} % {契丹武士要来中原夺武林秘笈。—— 慕容博}}
% \subsection{调整定理的样式}
% 字体、间距、证毕符号这些写在 \option{thm} 族里：
% \begin{latex}
% \hitsetup{
%   thm = {
%     body-font  = {\songti\rmfamily},
%     head-font  = {\heiti\rmfamily},
%     qed        = {$\square$},
%     head-punct = {},
%   },
% }
% \end{latex}
%
% 键名与 \pkg{thmtools} 的 \cs{declaretheoremstyle} 一致，只是改成了连字符分词：
%
% \begin{longtblr}[caption={\option{thm} 族的键}, label={tab:thm}]
%   {colspec={l l}}
% \toprule
% 键 & 管什么 \\ \midrule
% \option{body-font}   & 定理正文字体 \\
% \option{head-font}   & 定理标题字体 \\
% \option{note-font}   & 注记字体，留空则跟随标题 \\
% \option{note-braces} & 注记两侧的括号，两个参数 \\
% \option{head-punct}  & 标题末尾的标点 \\
% \option{head-space}  & 标题与正文之间的间距 \\
% \option{head-indent} & 标题的缩进 \\
% \option{qed}         & 证毕符号 \\
% \option{space-above} & 定理上方留白 \\
% \option{space-below} & 定理下方留白 \\
% \option{within}      & 编号跟随哪一级计数器 \\
% \option{list}        & 声明一批定理环境，写成\meta{环境名}\texttt{=}\meta{标题} \\
% \option{star}        & 同上，但声明的是不编号的 \\
% \bottomrule
% \end{longtblr}
%
% 设置只对之后声明的环境生效，所以要写在 \cs{newtheorem} 之前，
% 也在同一次 \option{thm} 里写在 \option{list} 之前。
%
% \subsection{\pkg{ntheorem} 的写法}
% v3.2a 之前定理环境由 \pkg{ntheorem} 实现。\pkg{ntheorem} 的 \cs{@endtrivlist}
% 与 \pkg{enumitem} 冲突，\env{proof} 结尾的证毕符号会丢（\issue[292]），定理块
% 末尾还会多出一段空白（\issue[294]），所以改用了 \pkg{amsthm} 与 \pkg{thmtools}。
%
% 下面这些 \pkg{ntheorem} 专有命令仍然可用，编译时会在日志里提示新写法，
% v4（预计~2027~年发布）移除：
%
% \begin{longtblr}[caption={\pkg{ntheorem} 旧命令与新写法对照}, label={tab:ntheorem}]
%   {colspec={l l}}
% \toprule
% 旧命令 & 新写法 \\ \midrule
% \cs{theorembodyfont}   & \option{thm} = \{ \option{body-font} \} \\
% \cs{theoremheaderfont} & \option{thm} = \{ \option{head-font} \} \\
% \cs{theoremsymbol}     & \option{thm} = \{ \option{qed} \} \\
% \cs{theoremseparator}  & \option{thm} = \{ \option{head-punct} \} \\
% \cs{thref}             & \cs{autoref} \\
% \cs{listtheorems}      & \cs{listoftheorems} \\
% \cs{theoremlisttype}   & \cs{listoftheorems} 的可选参数 \\
% \cs{theoremprework}    & \cs{declaretheorem} 的钩子 \\
% \cs{theorempostwork}   & \cs{declaretheorem} 的钩子 \\
% \cs{theoremclass}      & \cs{declaretheorem} 的 \option{style} 键 \\
% \cs{theoremnumbering}  & \cs{newtheorem} 之后重定义 \cs{the}\meta{环境名} \\
% \bottomrule
% \end{longtblr}
%
% \cs{theoremstyle} 现在是 \pkg{amsthm} 的命令，取值也换了一套，
% 原先 \pkg{ntheorem} 的 \texttt{break}、\texttt{change} 这些样式名不再认。
% 长度寄存器 \cs{theorempreskipamount} 与 \cs{theorempostskipamount} 也没有了，
% 请改用 \option{space-above} 与 \option{space-below}。
%
% \DescribeMacro{\quan}
% \DescribeMacro{\Quan}
% \DescribeMacro{\quansplit}
% \subsection{带圈数字}
% 脚注号用的带圈数字由 \cs{quansplit}\marg{数字} 生成，也可以直接调用。当前宋体有
% 现成的带圈字符（\texttt{U+2460} 起，多数字体到 10 或 20）就用字符，字形比画出来的
% 规整；超出范围才用 \pkg{circledsteps} 画一个，并按字符版的尺寸缩放，混在一行里
% 看不出接缝。画的时候数字取当前宋体，与正文一致。
%
% \cs{quan}\oarg{十位}\marg{个位} 与 \cs{Quan}\marg{数字} 是同一个东西的两种写法，
% 压缩比例自动算，不必自己分位数。
%
% \subsection{算法}
% 我工算法不在规范中要求且一千个评审老师有一千个算法格式喜好。详见
% \href{https://github.com/PlutoThesis/PlutoThesis#%E6%B2%A1%E6%9C%89%E6%98%8E%E7%A1%AE%E8%A6%81%E6%B1%82%E7%9A%84%E6%A0%BC%E5%BC%8F}{PlutoThesis}
% 中的各个实验室算法喜好举例。在此多说无益。
% \subsection{引用参考文献}
% \DescribeMacro{\inlinecite}
% 学校要求的参考文献引用有两种模式：（1）上标模式。比如``同样的工作有很
% 多$^{[1,2]}$\ldots''。（2）正文模式。比如``文[3] 中详细说明了\ldots''。其中上标
% 模式使用远比正文模式频繁，所以为了符合使用习惯，上标模式仍然用常规
% 的 \cs{cite}\marg{key}，而 \cs{inlinecite}\marg{key} 则用来生成正文模式。
%
% 关于参考文献模板推荐使用 \hologo{BibTeX}，关于中文参考文献需要额外增加一个 Entry:
% \texttt{language}，将其设置为 \texttt{zh} 用来指示此参考文献为中文，以
% 便 \file{hithesis.bst} 处理。如：
% \begin{latex}
% @INPROCEEDINGS{cnproceed,
%   author    = {王重阳 and 黄药师 and 欧阳峰 and 洪七公 and 段皇帝},
%   title     = {武林高手从入门到精通},
%   booktitle = {第~$N$~次华山论剑},
%   year      = 2006,
%   address   = {西安, 中国},
%   language      = "zh",
% }
%
% @ARTICLE{cnarticle,
%   AUTHOR  = "贾宝玉 and 林黛玉 and 薛宝钗 and 贾探春",
%   TITLE   = "论刘姥姥食量大如牛之现实意义",
%   JOURNAL = "红楼梦杂谈",
%   PAGES   = "260--266",
%   VOLUME  = "224",
%   YEAR    = "1800",
%   LANGUAGE    = "zh",
% }
% \end{latex}
%
% 注意如果不需要引用参考文献，请删除 \file{final.tex} 中 \cs{bibliography} 开头的两行，
% 以避免可能的编译错误。
%
% \subsection{空行}
% \DescribeMacro{\hitenter}
% 在 Word 里空一行靠的是单独按一次回车。\cs{hitenter} 模拟这个动作：产生一段
% 恰等于“空段自动行高”的竖直距离，不必手写任何长度值。Word 的规则是段落标记
% 字体（Times）的自然行高乘上当前段落的行距，段落贴网格时再凑成整格；空段继承
% 上一段的段落属性，所以行距与贴格默认都跟当前段落走，正文里就是 1.25 倍不贴格。
% 可选参数是回车几次：
% \begin{latex}
% 上一段。
% \hitenter[2]
% 隔了两个空行的下一段。
% \end{latex}
% 也收键值：\option{lines}、\option{font}（\option{size}、\option{asian-text-font}、
% \option{latin-text-font}）、\option{paragraph}（\option{spacing} 里的 \option{line-spacing} 与
% \option{at}，以及 \option{snap-to-grid}），文法与 \option{styles} 族相同。
% 空行恰好落在换页处时默认随页首丢弃（页首空行在论文版式里没有意义）；
% 想保留 Word 空段占据页首一行的原义，用星号版 \cs{hitenter*}：每个空段是一只
% 不可弃的空盒，页底装得下几行装几行，剩下的落到新页页首照常占位。抄 Word 稿子
% 时用星号版。图表前后那一行是类给的，空段另算，两者相加。
% 模板自己排摘要与关键词之间的空行用的就是这条命令。
%
% \subsection{列表环境}
% \DescribeEnv{itemize}
% \DescribeEnv{enumerate}
% \DescribeEnv{description}
% 为了适合中文习惯，模板将这三个常用的列表环境用 \pkg{enumitem} 进行了纵向间距压
% 缩。一方面清除了多余空间，另一方面用户可以自己指定列表环境的样式（如标签符号，
% 缩进等）。细节请参看 \pkg{enumitem} 文档，此处不再赘述。
%
% \subsection{行内代码}
% \DescribeOption{verb-space}
% \cs{verb} 里的一串照 Word 里能一字一字映射的写法排：Consolas 五号（等宽字体按 7/8
% 缩放，小四正文里正好五号；没装 Consolas 就是 Courier New），空格印成一格宽的记号
% （\verb|a b| 印出来是 \verb|a b|），记号后可断行；超过 16 个字符的串在 \texttt{/}、
% \texttt{.}、\texttt{-} 前面可断，分隔符落到下一行行首，头尾四个字符与扩展名不断；
% 连字符后只有 \option{compatibility-mode}\texttt{=word-2013} 可断；一段自己起行还装不下时按字形硬断。对应的 Word 写法是 Consolas 五号的 run、
% 空格敲 U+02FD 加 U+200B、长串分隔符前敲 U+200B，两边排出来的行相同。
%
% \option{style}\texttt{=\{verb-space=cell\}} 把记号换成不可见的不断行空格（Word 里
% 是 U+00A0 加 U+200B），仍占一格、仍可断。\cs{verb*} 与不带星号的画法相同。
% \env{verbatim} 环境里的空格不动。
% \section{后文}
%
% \subsection{结论}
% \DescribeEnv{conclusions}
% 结论之后为后文内容。
% \begin{latex}
% \begin{conclusions}
%
% 学位论文的结论作为论文正文的最后一章单独排写，但不加章标题序号。
%
% 结论应是作者在学位论文研究过程中所取得的创新性成果的概要总结，不能与摘要混为一
% 谈。博士学位论文结论应包括论文的主要结果、创新点、展望三部分，在结论中应概括论
% 文的核心观点，明确、客观地指出本研究内容的创新性成果（含新见解、新观点、方法创
% 新、技术创新、理论创新），并指出今后进一步在本研究方向进行研究工作的展望与设想
% 。对所取得的创新性成果应注意从定性和定量两方面给出科学、准确的评价，分（1）、
% （2）、（3）…条列出，宜用“提出了”、“建立了”等词叙述。
%
% \end{conclusions}
% \end{latex}
%
% \subsection{参考文献}
% 在后文中的参考文献是自动生成的，不需要用户干预，具体命令在 \file{final.tex} 中有
% 示例。
%
% \subsection{附录}
% \DescribeEnv{appendix}
% 所有的附录都插到这里来。因为附录会更改默认的 chapter 属性，而后面的{\heiti 个人简
%   历}又需要恢复，所以实现为环境可以保证全局的属性不受影响。
% \begin{latex}
% \begin{appendix}
% \input{final/backmatter/appA.tex}
% \end{appendix}
% \end{latex}
%
% \subsection{所发表文章}
% \DescribeEnv{publication}
% \DescribeEnv{pubgroup}
% 指南 1.6 说不同形式的成果“可单独分项分别列出”，分几组、每组叫什么由作者
% 定。条目本身推荐从 \file{.bib} 里出，格式与参考文献表同一条链：
% \begin{latex}
% \begin{publication}[final/backmatter/pub]
%   \begin{pubgroup}{发表的学术论文}
%     \item{pub:oxid}
%     \item{pub:sim}
%   \end{pubgroup}
%   \begin{pubgroup}{申请及已获得的专利}
%     \item{pub:patent}
%   \end{pubgroup}
% \end{publication}
% \end{latex}
% \env{publication} 的可选参数是成果条目那个库，\env{pubgroup} 的参数是这一组
% 的标题，组里逐条 \cs{item}\marg{文献键}。成果页与参考文献共用标点、西文题名大小写和
% 作者数上限等 \option{bib} 设置；编号的
% 列宽、左右对齐也与参考文献表同一处取值，不用另配一份。每组从 1 起编号。
%
% 组名自动带编号：研究生“（一）发表的学术论文”，本科“一、发表的学术论文”，
% 英文版“(1) ”，都是范例的写法。可选参数改它：\texttt{[]} 不编号，
% \texttt{[甲]} 之类就用字面的。组名段的段前段后 5 磅、1.25 倍行距、宋体
% 加粗（Word 那种描边合成，不换黑体）都照范例；条目标号后空两个宋体空格、
% 折行悬挂 2.3 字，与参考文献表那边（标号后不空、悬挂到标号宽）是范例自己
% 定的两个数，不用你管。
%
% \textbf{成果条目要单独一个库。}它们走自己的一次 \pkg{bibtex}，与参考文献表
% 彻底分开：主参考文献表从头到尾看不见这些条目。编译要多跑一句
% \texttt{bibtex}~\meta{jobname}\texttt{-pub}，仓库里的 \file{Makefile} 与
% \file{latexmkrc} 都已经带上，用 \pkg{latexmk} 或 \texttt{make} 的话不必自己管。
%
% \textbf{\pkg{biblatex} 一线（\option{bib}/\option{backend}\texttt{=biber}）
% 写法不同：库名写在导言区。}
% \begin{latex}
% \addbibresource{final/backmatter/pub.bib}   % 导言区
% ...
% \begin{publication}                          % 不写可选参数
%   \begin{pubgroup}{发表的学术论文}
%     \item{pub:oxid}
%   \end{pubgroup}
% \end{publication}
% \end{latex}
% 这一线不用多跑那一次 \pkg{bibtex}，条目由 \cs{fullcite} 取，但它只认已经
% \cs{addbibresource} 登记过的库，而 \cs{addbibresource} 只能在导言区用——%
% \pkg{biblatex} 在 \cs{begin\{document\}} 那一刻就把资源表交给 \pkg{biber} 了，
% 而 \env{publication} 是正文里的环境，那时早过了。在环境上写了库名会收到一条
% 提示叫你挪到导言区，不是静默忽略：静默的话每一条成果都排成空白，只在
% \file{.blg} 里留一句 \texttt{I~didn't~find~a~database~entry}，没人会去看。
%
% 不给那个库名就是老写法，正文随作者写，条目用 \env{publist} 逐条手敲：
% \begin{latex}
% \begin{publication}
% \noindent\textbf{（一）发表的学术论文}
% \begin{publist}
% \item	XXX，XXX. Static Oxidation Model of Al-Mg/C Dissipation Thermal Protection Materials[J]. Rare Metal Materials and Engineering, 2010, 39(Suppl. 1): 520-524.（SCI~收录，IDS号为~669JS，IF=0.16）
% \item XXX，XXX. 精密超声振动切削单晶铜的计算机仿真研究[J]. 系统仿真学报，2007，19（4）：738-741，753.（EI~收录号：20071310514841）
% \item XXX，XXX. 局部多孔质气体静压轴向轴承静态特性的数值求解[J]. 摩擦学学报，2007（1）：68-72.（EI~收录号：20071510544816）
% \item XXX，XXX. 硬脆光学晶体材料超精密切削理论研究综述[J]. 机械工程学报，2003，39（8）：15-22.（EI~收录号：2004088028875）
% \item XXX，XXX. 基于遗传算法的超精密切削加工表面粗糙度预测模型的参数辨识以及切削参数优化[J]. 机械工程学报，2005，41（11）：158-162.（EI~收录号：2006039650087）
% \item XXX，XXX. Discrete Sliding Mode Cintrok with Fuzzy Adaptive Reaching Law on 6-PEES Parallel Robot[C]. Intelligent System Design and Applications, Jinan, 2006: 649-652.（EI~收录号：20073210746529）
% \end{publist}
%
% \noindent\textbf{（二）申请及已获得的专利（无专利时此项不必列出）}
% \begin{publist}
% \item XXX，XXX. 一种温热外敷药制备方案：中国，88105607.3[P]. 1989-07-26.
% \end{publist}
%
% \noindent\textbf{（三）参与的科研项目及获奖情况}
% \begin{publist}
% \item	XXX，XXX. XX~气体静压轴承技术研究, XX~省自然科学基金项目.课题编号：XXXX.
% \item XXX，XXX. XX~静载下预应力混凝土房屋结构设计统一理论. 黑江省科学技术二等奖, 2007.
% \end{publist}
% %\vfill
% %\hangafter=1\hangindent=2em\noindent
% %\setlength{\parindent}{2em}
% \end{publication}
% \end{latex}
%
% \subsection{索引}
% \DescribeEnv{ceindex}
% 我工要求中英文双语索引。后文中的自动索引实际上不需要用户干预。
%\begin{latex}
% \begin{ceindex}
%   %如果想要手动加索引，注释掉以下这一样，用wordlist环境
% \printsubindex*
% \end{ceindex}
%\end{latex}
% 手工添加索引的方法不推荐，模板中将去除该功能。
% \subsection{原创性声明与使用权限}
% \DescribeOption{declarations}
% 这一页有两节：原创性声明与使用权限。整页叫 \option{declarations}，写在
% \option{struct} 里：
%\begin{latex}
% struct = { backmatter = {
%   declarations  = true,                  % 排出来
%   declarations  = {figures/scan.pdf},    % 贴这个 PDF，不排页码
%   declarations* = {figures/scan.pdf},    % 贴这个 PDF，排页码
% } }
%\end{latex}
% 签名与日期要手写，排版这边不用管；真要在编译时填死就改
% \option{term} 里的 \option{author-signature-field-label}、
% \option{supervisor-signature-field-label}、\option{date-field-label-front}、
% \option{date-blank}，最后那个只有这一页在用。
%
% \textbf{页内两节的文案分开配。} 原创性声明那一节的键带 \option{copyright}，
% 使用权限那一节的带 \option{authorization}：
%
% \begin{longtblr}[caption={\option{declarations} 部件的键}, label={tab:declarations}]
%   {colspec={X[2,l] X[1,l]}, width=\linewidth}
% \toprule
% 键 & 管哪一块 \\ \midrule
% \option{declarations-heading-zh}、\option{declarations-heading-en} & 整页标题 \\
% \option{statement-heading} & “学位论文原创性声明”那行小标题 \\
% \option{statement-text} & 原创性声明的正文 \\
% \option{authorization-heading} & “学位论文使用权限”那行小标题 \\
% \option{authorization-text} & 使用权限的正文 \\
% \bottomrule
% \end{longtblr}
%
% 本科那一页的对应键前面多一个 \option{bachelor-}，各校区文案不同，
% 默认取哈尔滨的，威海与深圳各自覆盖。
%
% \DescribeMacro{\authorization}
% \cs{authorization}\oarg{文件} 与 \cs{authorization*}\oarg{文件} 照旧能用，
% 但已经是旧写法：调用时日志里会提示改写成上面那个键，
% 兼容期到 v4（预计~2027~年发布）。
%
% \subsection{致谢声明}
% \DescribeEnv{acknowledgement}
% 把致谢做成一个环境更好一些，直接往里面写感谢的话就可以啦！
%
% \begin{latex}
% \begin{acknowledgement}
%   …
%   感谢\hit\LaTeX\ 论文模板\hithesis\ !
% \end{acknowledgement}
% \end{latex}
%
% \subsection{简历}
% \DescribeEnv{resume}
% 个人简历。
% 实际上，致谢和个人简历是自由发挥的地区，字体，文体，格式，内容，完全自己决定。
% \begin{latex}
% \begin{resume}
% XXXX~年~XX~月~XX~日出生于~XXXX。
%
% XXXX~年~XX~月考入~XX~大学~XX~院（系）XX~专业，XXXX~年~XX~月本科毕业并获得~XX~学学士学位。
%
% XXXX~年~XX~月------XXXX~年~XX~月在~XX~大学~XX~院（系）XX~学科学习并获得~XX~学硕士学位。
%
% XXXX~年~XX~月------XXXX~年~XX~月在~XX~大学~XX~院（系）XX~学科学习并获得~XX~学博士学位。
%
% 获奖情况：如获三好学生、优秀团干部、X~奖学金等（不含科研学术获奖）。
%
% 工作经历：
% \end{resume}
% \end{latex}
%
% \section{改写词条}
% \label{sec:const}
% 封面标签、章节名、声明正文这类固定词汇叫\emph{词条}，由模板给出默认值。各院系
% 的要求不尽相同，需要改写时在导言区用 \option{term} 族设置：
% \begin{latex}
% \hitsetup{
%   term = {
%     zh = { supervisor-field-label = {指导教师}, keywords-separator = {；} },
%   },
% }
% \end{latex}
% \emph{这一族原先叫 \option{const}。}改名的理由是族名该说清“这一族键是什么”：
% 每个键都是一条词条，取出来的宏一直叫 \cs{hit@term@}\meta{名}，而
% \texttt{const}（常量）只说了它不会变，没说是哪一类东西。旧族名照旧生效，
% 写了会在日志里提示一次；只有斜杠路径那种写法（\texttt{const/zh/\ldots}）
% 不认旧名，要写成 \texttt{term/zh/\ldots}。
%
% 中英文成对的词汇同样收在 \option{zh} 与 \option{en} 分组里，写法见
% 第~\ref{sec:setup-language}~节。带后缀的本名与内部命令一一对应：
% \option{supervisor-field-label-zh} 对应 \cs{hit@supervisor@field@label@zh}，
% 名字里的连字符对应内部命令里的 \texttt{@}。
%
% \emph{键名怎么拼。} 一个键名是“桶名 + 词条名 + 各档取值”。桶名在最前，
% 各档取值一律排在后面，次序固定：
%
% \begin{quote}
% \meta{词条名}\oarg{-材料}\oarg{-校区}\oarg{-学位}\oarg{-学位类别}\oarg{-学科门类}\oarg{-语言}
% \end{quote}
%
% 省掉哪一段，就是那条轴各档共用。\option{declarations-statement} 是研究生与本科
% 共用的那一份，\option{declarations-statement-bachelor} 是本科专有的那一份；
% \option{school-bottom-mark-shenzhen-master} 同时钉住校区与学位。
%
% 六根轴与六个类选项一一对应，取值也一样：材料是 \option{stage}
% （\texttt{final}／\texttt{proposal}／\texttt{interim}），校区是 \option{campus}，
% 学位是 \option{degree-level}，学位类别是 \option{degree-type}，学科门类是
% \option{category}，语言是 \option{lang}。
%
% \emph{档位值前面加 \texttt{!} 就是取反。} \option{achievements-heading-!final}
% 读作“不是终稿的那几档”，也就是开题与中期共用的那一条。终稿是这根轴的基准，
% 不为它另立名字——与条件里 \option{layout}\texttt{=!v2-layout} 是同一套写法。
%
% 比具体程度时，正着写的一段压得住取反的那一段：同时写了
% \option{achievements-heading-proposal} 与 \option{achievements-heading-!final}，
% 开题那一档取前者。语种段不许取反，写了当场报错。
%
% \emph{桶是页。} 同一个字段身份在不同的页上常常是不同的东西：报告封面的“院（系）”
% 与内封的“所在单位”都是 \option{affiliation}，挤在一层里只能靠前缀区分，越加越长，
% 也看不出哪一条只有某一页用得上。所以键名前面带一段桶名，一页一个桶，跨页共用的
% 进 \texttt{base}。桶名共十个：\texttt{base}、\texttt{cover}、\texttt{titlepage}、
% \texttt{abstract}、\texttt{toc}、\texttt{declarations}、\texttt{resolution}、
% \texttt{backmatter}、\texttt{header}、\texttt{caption}。
%
% 写的时候桶与语言各是一层，不必自己拼：
%
% \begin{latex}
% \hitsetup{ term = {
%   zh = { titlepage = { supervisor = {导师} } },   % titlepage-supervisor-zh
%   declarations = { authorization = {使用授权书} },  % declarations-authorization
% } }
% \end{latex}
%
% 三种写法等价：分层写、拼好写（\option{declarations-authorization}）、斜杠路径写
% （\option{term/zh/titlepage/supervisor}）。拼错桶名或词条名会收到一条警告，
% 里面写着拼出来的完整键名。
% 模板按当前文档的校区、学位、语种挑段数最多的一条，绑到不带任何档位的名字上，
% 所以排版代码里写 \cs{hit@statement@heading} 拿到的就是本档该用的值。
%
% 用户改写时两种名字都收：写不带档位的名字覆盖当前这一档，写带档位的名字只覆盖
% 那一档。
%
% v3.2a 开发中有 27 个键把档位写在名字前头（\option{bachelor-statement-heading}
% 这样），已经统一改成后缀。这些键都是 v3.2a 新加的，v3.1e 没有对应的写法。
%
% \subsection{默认值都在 \file{hithesis.cfg} 里}
% \label{sec:cfg}
% 模板给的默认值全在 \file{hithesis.cfg} 里，一行一个键，两档都读这一份。想知道
% 某个词、某个尺寸的默认值是什么，在里面搜键名就能看到；要改，照第
% ~\ref{sec:keyfamily}~节在导言区用 \cs{hitsetup} 覆盖，不必改这个文件。
%
% 它由 \file{hithesis.dtx} 生成，跟 \file{hithesis.cls} 一起装进
% \file{tex/latex/hithesis/}。里面每一块的形式是
% \cs{hit@presetup}\oarg{条件}\marg{键值表}：条件的写法与
% \cs{hitsetup} 的可选参数完全一样，见第~\ref{sec:setup-condition}~节。
% 不带条件的那些块无条件生效。
%
% 文件分三段。头一段不带段条件，两档都吃；接下来
% \cs{hit@presetup@scope}\texttt{\{stage!=final\}} 起是终稿那一段，
% \cs{hit@presetup@scope}\texttt{\{stage!=!!final\}} 起是开题与中期那一段。
% 段条件与每一行自己的条件各判各的，都成立才生效。
%
% v3.1e 是 \file{hithesisbook.cfg} 与 \file{hithesisart.cfg}，v3.2a 开发中一度
% 并进类文件，又按类名拆成两份，最后合成这一份。
%
% \subsection{编号之后的间距}
% \label{sec:after-number}
% \DescribeOption{after-chapter-number}
% \DescribeOption{after-section-number}
% \DescribeOption{after-subsection-number}
% \DescribeOption{after-subsubsection-number}
% \DescribeOption{after-heading-number}
% \DescribeOption{after-figure-caption-number}
% \DescribeOption{after-table-caption-number}
% “第 1 章”与章标题之间、“图 2-1”与图题之间排多宽，原先写死在类里。这批常量
% 把它们放出来，同样收在 \option{term} 族，同样分中英：
% \begin{latex}
% \hitsetup{term={zh={after-chapter-number={\quad}}}}
% \end{latex}
%
% \begin{longtblr}[caption={编号之后的间距词条}, label={tab:after-number}]
%   {colspec={l l}}
% \toprule
% 常量 & 管哪个编号之后 \\ \midrule
% \option{after-chapter-number} & 章，只有终稿有 \\
% \option{after-section-number} & 节 \\
% \option{after-subsection-number} & 小节 \\
% \option{after-subsubsection-number} & 小小节 \\
% \option{after-heading-number} & 上面四个一起设 \\
% \option{after-figure-caption-number} & 图题 \\
% \option{after-table-caption-number} & 表题 \\
% \bottomrule
% \end{longtblr}
%
% \option{after-heading-number} 等于把四级逐一写一遍，所以后面再单设某一级仍然
% 管用，先铺后改的写法是有效的：
% \begin{latex}
% \hitsetup{term={zh={after-heading-number={\enspace},
%                      after-chapter-number={\quad}}}}
% \end{latex}
%
% 取值是排版内容，写 \cs{quad}、\cs{enspace}、\cs{hskip}\cs{ccwd} 这类间距最常见，
% 写别的也行。
%
% 不设或设成空值时按模板原有的规矩排，那套规矩不是一个固定值：新版面的章节标题照
% 原件在题名前敲两个半角空格，合起来一个字宽，英文学位论文的章号之后换行，人文社科类
% 的节与条题序后不空；旧版面章号之后本部硕士与深圳本科是一个全角、其余是半个；图表
% 题号之后本科是一个汉字宽、其余是 \cs{enskip}，报告还要看开题还是中期。所以这批
% 常量的默认值是空，空的含义是“照旧”，不是“什么都不排”。要真的不排，写一对空的
% 花括号是不够的，得写 \cs{hspace}\marg{0pt} 这种明确为零的值。
%
% \section{捐助}
% \changes{v1.0.01}{2017/08/27}{添加了捐助、矢量化本科论文模板的图片logo}
% 各位刀客和大侠如用的嗨，要解囊相助，请支付宝赞助
% 参照图~\ref{zfb}~中提示操作（二维码被矢量化后之后去
% 除了头像等冗余无用的部分～）。
% \begin{figure}[h]
% \centering
\includegraphics[width=0.5\textwidth]{zfb}
% \caption{支付宝}
% \label{zfb}
% \end{figure}
%
